% Korean Noether paper 22 — fresh producer translation, unchecked.
% Initial translation authority: ED0005 lines 12702--13533.
% Current scope-rebase authority: ED0011 / NOETH-DE-AUTH-v047-20260808.
% Authority SHA-256: 1A44F967B29972E8F99E5C323A479162AD82A23FC457395915A4BB9DDF51AD41.
% Normalized source span: 80,281 bytes; SHA-256 759251216C47DD8CE2FE87AA92319C61CAEE976325FF309D0197A8301B557914.
% Translation only; independent source, Korean-language, render, and certification checks remain external.

\providecommand{\Amod}{\mathfrak A}
\providecommand{\Bmod}{\mathfrak B}
\providecommand{\Gmod}{\mathfrak G}
\providecommand{\Mmod}{\mathfrak M}
\providecommand{\Nmod}{\mathfrak N}
\providecommand{\Rmod}{\mathfrak R}
\providecommand{\Smod}{\mathfrak S}
\providecommand{\Tmod}{\mathfrak T}
\providecommand{\mideal}{\mathfrak m}
\providecommand{\nideal}{\mathfrak n}
\providecommand{\gideal}{\mathfrak g}
\providecommand{\hideal}{\mathfrak h}
\providecommand{\jideal}{\mathfrak j}
\providecommand{\tideal}{\mathfrak t}
\providecommand{\aideal}{\mathfrak a}
\providecommand{\bideal}{\mathfrak b}
\providecommand{\bb}{\mathfrak b}
\providecommand{\cc}{\mathfrak c}
\providecommand{\zero}[1]{\equiv0\pmod{#1}}
\providecommand{\Norm}{\operatorname N}

\section*{22. 다항식 아이디얼과 종결식 이론}

\begin{center}
\emph{Math. Ann. 88 (1923), 53--79쪽}\\[0.9em]
쿠르트 헨첼(\(\dagger\)) 지음\\[0.6em]
괴팅겐의 에미 뇌터 편집\srcfn{1)}{쿠르트 헨첼은 1914년 10월 이래 딕스뮈드 전선에서 실종되었으며 사망자로 보아야 한다. 이 글은 그가 1914년 여름 에를랑겐에서 E. Fischer의 지도 아래 제출한 학위논문 \glqq{}\emph{Zur Theorie der Formenmoduln und Resultanten}\grqq{}의 가장 본질적인 부분을 완전히 자유롭게 다시 쓴 것이다. 전적으로 독창적인 착상으로 쓴 그 학위논문은 빈틈없이 구성되었지만, 보조정리가 보조정리를 잇고 모든 개념이 네댓 개 지표를 가진 공식으로만 서술되며 본문이 거의 전혀 없다. 그래서 이해하기가 극도로 어렵다. 본인은 계획한 개작을 완수하지 못했다. 나는 이 연구를 순전히 개념적인 형태로 다시 제시한다. 이로써 기본착상이 한결같이 헨첼에게서 비롯된 증명들이 크게 단순해지고, 바라건대 이 연구의 아름다움도 드러날 것이다. 주어진 기저에서 유한개의 단계로 나타나는 함수들을 실제로 구성하는 문제를 다룬 학위논문의 부분은 별도 출판을 위해 남겨 둔다. (E. N.)}
\end{center}

이하에서는 \emph{다항식 아이디얼에 속하는 모든 다항식의 공통영점을 정한다는 소거론의 일반문제}를 다룬다. 이는 그 아이디얼의 기저를 이루는 임의의 유한개 다항식의 영점을 정하는 것과 같다. 동시에 나타나는 중복도에 대한 개념적 해석도 얻을 것이다. 다항식의 계수는 임의의 체에 속해도 되며, 영점은 그 체에서 유도되는 대수적으로 닫힌 체에 속한다. 또한 일반성을 잃지 않고 아이디얼을 변환된 아이디얼(제3절)이라고 가정할 수 있다. 그러면 본질적인 결과는 다음과 같다.

\emph{각 다항식 아이디얼 \(\mideal\)에 종결식형
\[
R_{\mideal}
=R^{(1)}(x_1,\ldots,x_n)\cdots
R^{(i)}(x_i,\ldots,x_n)\cdots R^{(n)}(x_n)
\zero{\mideal}
\]
을 대응시킬 수 있으며, \(R_{\mideal}\)은 \(\mideal\)의 모든 영점에서, 그리고 오직 그 영점들에서만 사라진다. \(\nideal\)이 \(\mideal\)의 인자이고 종결식형 \(R_{\nideal},R_{\mideal}\)이 같으면 아이디얼 \(\nideal,\mideal\)도 같다.}

주정리의 둘째 부분은 종결식형이 각 영점을 특징적인 중복도와 함께 준다는 것을 보여 준다. 이는 한 아이디얼이 다른 아이디얼의 인자인 두 대수적 수체 아이디얼이 그 \emph{노름}이 같을 때 그리고 그때에만 같다는 사실의 유사물이며, 두 정리의 내적 근거도 같다.

실제로 각 다항식을 \(x_1,\ldots,x_{i-1}\)의 거듭제곱곱들에 관한 선형형식으로 보고 \(x_{i+1},\ldots,x_n\)을 계수영역에 첨가하면, \(\mideal\)을 선형형식 모듈 \(\mathfrak M_{i-1}\)로 볼 수 있다. 종결식인자 \(R^{(i)}\)는 그때 대응하는 기초모듈(제1절) \(\Gmod_{i-1}\)의 \(\mathfrak M_{i-1}\)에 대한 \emph{노름}과 같아진다. 따라서 \emph{\(R^{(i)}\)의 한 인자의 차수와 지수의 곱으로 주어지는 중복도를 선형독립인 잉여류의 수로 해석}할 수 있다\srcfn{2)}{이 마지막 결과들의 진정한 의미는 아이디얼의 으뜸아이디얼 분해를 함께 고려할 때 드러난다. 한 으뜸아이디얼에 대응하는 중복도는 그 아이디얼에 대한 보충부분의 선형독립인 잉여류 수로 정의된다. 이는 다른 곳에서 다루겠다. (E. N.)}. 지금까지 이 사실은 부정계수에 대하여 종결식을 정의할 수 있는 특수한 경우에만 알려져 있었다\srcfn{3)}{Macaulay, \emph{The algebraic theory of modular systems}, Cambridge Tracts 19 (1916), 67번; Lasker, \emph{Zur Theorie der Moduln und Ideale}, \emph{Math. Ann.} 60 (1905), 20--116쪽, 특히 98쪽. 이 경우 종결식과 종결식형이 일치한다는 것은 주 1에서 말한 아직 출판되지 않은 헨첼의 정리들에 들어 있다. (E. N.)}.

이 결과들을 유도하기 위해 제1절과 제2절에서는 계수가 유리정수 또는 한 부정원의 다항식인 선형형식 모듈에 관한 정리들을 제시한다. 제3절에서 도입하는 다항식 아이디얼과의 연결은 제4절의 동형정리가 맡는다. 제5절은 위에서 말한 주정리의 둘째 부분을, 제7절은 영점에 관한 정리를 준다. 그에 앞선 제6절에서는 일반 소거론을 전개한다. 여기서는 새 부정원들을 도입하여 레졸벤트가 이 부정원들의 일차형식들로 분해됨도 증명한다. 이로써 여기서 제시하는 소거법 안에서 크로네커의 한 주장에 엄밀한 증명을 준다\srcfn{4)}{König, \emph{Algebraische Größen} (Leipzig: Teubner, 1903), V, \S{}4가 크로네커 소거론에 대하여 시도한 증명은 성공하지 못한 것으로 알려져 있다. König이 크로네커 소거론에 도입한 중복도들에 대해서도 헨첼 주정리의 둘째 부분이 성립하지 않는다는 사실은 Macaulay, 앞의 책 23쪽에 나와 있다. 거기서는 크로네커 소거론이 으뜸아이디얼 분해와 대응하지 않는다는 것도 보인다. (E. N.)}.

\section*{\S{}1. 선형형식의 모듈}

앞의 두 절에서 정수적 원소 \(a,b,c,\ldots\)란 유리정수 또는 추상적으로 정의된 임의의 체 \(P\)를 계수체로 갖는 한 부정원의 다항식을 뜻한다. 선형형식 \(a(\xi),b(\xi),\ldots\)란 정수적 원소를 계수로 갖는 유한개의 부정원 \(\xi_1,\ldots,\xi_k\)의 선형형식을 뜻한다.

\emph{선형형식의 모듈 \(\Amod\)}은 다음 조건으로 정의한다. \(\Amod\)가 \(a(\xi),b(\xi)\)를 포함하면 차 \(a(\xi)-b(\xi)\)도 포함한다. \(a(\xi)\)를 포함하면 임의의 정수적 원소 \(c\)에 대하여 \(c\,a(\xi)\)도 포함한다.

\(\Amod\)의 \emph{랭크}는 통상처럼 \(\Amod\)에 속하는 선형독립인 선형형식의 최대수로 정의한다. \(\Amod\)의 \(\lambda\)번째 \emph{행렬식인자} \(d_\lambda\)는 \(\Amod\)에서 얻는 모든 \(\lambda\)차 행렬식의 최대공약원이다. 즉 \(\Amod\)의 선형형식 \(\lambda\)개씩의 계수행렬에서 만들 수 있는 모든 \(\lambda\)차 행렬식을 취한다. \(\varrho\)가 \(\Amod\)의 랭크여서 \(d_1,\ldots,d_\varrho\)가 영이 아닌 행렬식인자들이라면, \(\Amod\)의 \emph{초등인자}는
\[
e_1=d_1,\quad e_2=\frac{d_2}{d_1},\ldots,\quad
e_\varrho=\frac{d_\varrho}{d_{\varrho-1}},\qquad
e_{\varrho+i}=0
\]
으로 정의한다. 잘 알려져 있듯 \(\Amod\)은 정확히 \(\varrho\)개의 선형형식 \(a_1(\xi),\ldots,a_\varrho(\xi)\)로 이루어진 모듈기저를 갖는 유한모듈이다. \(\Amod\)의 모든 선형형식은 이 기저를 정수적 원소의 계수로 선형표현할 수 있다.
\[
\Amod=(a_1(\xi),\ldots,a_\varrho(\xi)).
\]
\(A\)가 이 선형형식들의 계수행렬이면 \(\Amod\)의 행렬식인자와 초등인자는 \(A\)의 것들과 같다. 먼저 여기서 \emph{각 초등인자 \(e_i\)가 다음 초등인자 \(e_{i+1}\)을 나눈다}는 것이 따라온다. 초등인자론의 행렬방정식
\[
PAQ=
\begin{pmatrix}
e_1&&&0\\
&e_2&&\\
&&\ddots&\\
0&&&e_\varrho
\end{pmatrix}
\]
에서 단모듈러변환 \(\xi=Q(\eta)\)로 새 부정원 \(\eta_1,\ldots,\eta_k\)를 도입하면 모듈등식
\begin{equation}
\Amod=(e_1\eta_1,e_2\eta_2,\ldots,e_\varrho\eta_\varrho)
\tag{1}
\end{equation}
을 얻는다.

이하에서 가장 중요한 개념은 기초모듈의 개념이다\srcfn{5)}{Steinitz, \emph{Rechteckige Systeme und Moduln in algebraischen Zahlkörpern II}, \emph{Math. Ann.} 72 (1912), 297--345쪽, 36번 참조. 다른 접점들도 있지만, 적어도 헨첼은 더 단순한 자신의 경우에서 슈타이니츠와 독립적으로 식 (4)의 형태로 이 개념에 도달한 듯하다. (E. N.)}.

\medskip
\noindent\textbf{정의 I.} \emph{같은 랭크의 진인자를 갖지 않는 모듈 \(\Gmod\)을 기초모듈이라 한다}\srcfn{6)}{여기서 인자는 모듈의 의미이다. \(\Amod\)의 모든 원소가 \(\Bmod\)에 들어갈 때 \(\Amod\)은 \(\Bmod\)로 나누어진다고 쓰며, \(\Amod\zero{\Bmod}\)이다. \(\Bmod\)가 \(\Amod\)에 없는 원소를 포함하면 진인자라고 한다.}.

따라서 \(\Gmod\)이 기초모듈인 것은
\begin{equation}
c\,g(\xi)\zero{\Gmod},\qquad c\ne0
\quad\Longrightarrow\quad g(\xi)\zero{\Gmod}
\tag{2}
\end{equation}
가 언제나 성립할 때 그리고 그때뿐이다.

임의의 모듈과 기초모듈의 관계는 다음 정리가 준다.

\medskip
\noindent\textbf{정리 I.} \emph{모든 모듈 \(\Amod\)은 같은 랭크의 유일한 기초모듈 \(\Gmod\)로 나누어진다. \(\Gmod\)은 \(\Amod\)의 같은 랭크를 가진 최소인자로 정의되며
\begin{equation}
\Gmod=(\eta_1,\ldots,\eta_\varrho)
\tag{3}
\end{equation}
로 표현된다. \(\Gmod\)을 \(\Amod\)의 기초모듈이라 부른다.}

\(\Gmod\)이 같은 랭크의 최소인자라는 조건은 다음처럼 표현된다. \(\Gmod\)은 관계
\begin{equation}
b\,g(\xi)\zero{\Amod},\qquad b\ne0
\tag{4}
\end{equation}
를 만족하는 선형형식 \(g(\xi)\)만을, 그리고 그 모두를 포함한다. 여기서 먼저 \(\Gmod\)이 모듈임을 알 수 있다. 실제로
\[
b_1g_1(\xi)\zero{\Amod},\quad b_1\ne0,\qquad
b_2g_2(\xi)\zero{\Amod},\quad b_2\ne0
\]
이면
\[
b_1b_2\bigl(g_1(\xi)-g_2(\xi)\bigr)\zero{\Amod},
\qquad b_1b_2\ne0
\]
이고, 물론 \(b_1c\,g_1(\xi)\zero{\Amod}\)도 성립한다. \(\Gmod\)은 또한 기초모듈이다. 실제로
\[
c\,g(\xi)\zero{\Gmod},\qquad c\ne0
\]
이면 (4)에 따라
\[
bc\,g(\xi)\zero{\Amod},\qquad bc\ne0
\]
이고, 다시 (4)에 따라 \(g(\xi)\zero{\Gmod}\)이다.

조건들을 만족하면서 같은 랭크의 최소인자 \(\Gmod\)과 다른 기초모듈 \(\Gmod_1\)은 존재하지 않는다. \(\Gmod_1\)은 \(\Gmod\)로 나누어져야 하지만 기초모듈로서 같은 랭크의 진인자를 갖지 않으므로 \(\Gmod\)과 같다. 끝으로 (3)의 오른쪽이 나타내는 모듈은 기초모듈이다. 그 진인자는 또 다른 부정원 \(\eta\)를 포함해야 하므로 더 높은 랭크를 갖기 때문이다. 따라서 (3)은 \(\Amod\)의 유일하게 정의된 기초모듈을 주며, 이로써 \emph{정리 I의 모든 부분이 증명되었다}.

이하에서 또 하나 중요한 개념은 두 모듈의 몫이라는 데데킨트의 개념이다\srcfn{7)}{데데킨트는 곱셈도 정의되는 수모듈을 다루므로 그의 몫은 여기서 정의하는 몫과 일치하지 않는다. (E. N.)}.

\medskip
\noindent\textbf{정의 II.} \emph{두 선형형식 모듈의 몫 \(\cc=\Amod/\Bmod\)은 조건
\[
c\,\Bmod\zero{\Amod}
\]
을 만족하는 정수적 원소 \(c\) 전체로 정의한다. \(\cc\)는 정수적 원소들의 아이디얼, 따라서 주아이디얼이다. 이에 대응하여 선형형식 모듈 \(\Amod\)을 정수적 원소의 아이디얼 \(\bb\)로 나눈 몫 \(\mathfrak C=\Amod/\bb\)는 조건
\[
c(\xi)\,\bb\zero{\Amod}
\]
을 만족하는 선형형식 \(c(\xi)\) 전체로 정의한다. \(\mathfrak C\)는 다시 선형형식 모듈이며, \(\bb\)가 영아이디얼이 아니면 \(\Amod\)과 같은 랭크를 갖는다.}

정의에 따라 각 몫이 모듈의 성질을 갖는다는 것은 분명하다. 한편 정수적 원소들로 이루어지고 모듈의 성질을 갖는 계는 아이디얼이다.

몫 개념은 기초모듈과 가장 높은 초등인자 사이에 다음 관계를 준다.

\medskip
\noindent\textbf{정리 II.} \emph{\(\Amod\)의 기초모듈과 가장 높은 초등인자 사이에는 역관계
\begin{equation}
(e_\varrho)=\Amod/\Gmod,\qquad
\Gmod=\Amod/(e_\varrho)
\tag{5}
\end{equation}
가 성립한다. 여기서 \((e_\varrho)\)는 \(e_\varrho\)가 생성하는 주아이디얼이다.}

모든 초등인자가 \(e_\varrho\)를 나누므로 (1), (3)에 따라
\begin{equation}
e_\varrho\Gmod\zero{\Amod},
\qquad\text{따라서}\qquad
(e_\varrho)\Gmod\zero{\Amod}
\tag{6}
\end{equation}
이다. 그러므로 \((e_\varrho)\)는 몫 \(\Amod/\Gmod\)으로 나누어진다. 역도 성립한다. 실제로
\[
b\Gmod\zero{\Amod}
\quad\Longrightarrow\quad
b\eta_\varrho\zero{\Amod}
\quad\Longrightarrow\quad
b\zero{(e_\varrho)}.
\]
이로써 (5)의 첫째 공식이 증명된다\srcfn{8)}{\(\Amod_0=\Amod,\ldots,\Amod_\lambda=(\Amod,\eta_\varrho,\ldots,\eta_{\varrho-\lambda+1})\)로 두자. 여기서 \(\Amod_i\)의 가장 높은 초등인자는 \(e_{\varrho-i}\)이다. 같은 추론으로
\[
(e_\varrho)=\Amod_0/\Amod_1,\ldots,\quad
(e_{\varrho-\lambda})=\Amod_\lambda/\Amod_{\lambda+1},\ldots,\quad
(e_1)=\Amod_{\varrho-1}/\Amod_\varrho
\]
를 얻는다. 더 일반적으로 \(\zeta_\lambda=b_{\lambda1}\eta_1+\cdots+b_{\lambda\lambda}\eta_\lambda\)로 두고 \((b_{\lambda\lambda},e_\lambda)=1\)이라 하며
\[
\Bmod_0=\Amod,\ldots,\qquad
\Bmod_\lambda=(\Amod,\zeta_\varrho,\ldots,\zeta_{\varrho-\lambda+1})
\]
로 두자. 그러면 \(\Bmod_\lambda\)의 가장 높은 초등인자도 \(e_{\varrho-\lambda}\)이므로
\[
(e_{\varrho-\lambda})=\Bmod_\lambda/\Bmod_{\lambda+1}
\]
이다.}. 또한 (6)은 \(\Gmod\)이 몫 \(\Amod/(e_\varrho)\)으로 나누어짐을 보인다. 이 몫은 \(\Amod\)의 같은 랭크를 가진 인자이므로, \(\Amod\)의 기초모듈 정의에 따라 거꾸로 \(\Gmod\)로도 나누어진다. 이로써 (5)의 둘째 공식도 증명되었다.

\(d\)가 \(\Amod\)에서 얻는 항등적으로 영이 아닌 임의의 \(\varrho\)차 행렬식이라 하자. 그러면 \(d\)는 \(\varrho\)번째 행렬식인자 \(d_\varrho\), 따라서 \(e_\varrho\)로 나누어진다. 그러므로 \(d\Gmod\zero{\Amod}\)이고 앞의 추론에서 \(\Gmod=\Amod/(d)\)가 따라온다. 뒤에서 중요한 점은 \(\Gmod\)의 이 몫표현을 초등인자로 돌아가지 않고도 증명할 수 있어서 더 일반적으로 성립한다는 사실이다.

\medskip
\noindent\textbf{정리 III.} \emph{정수적 원소를 여러 부정원의 다항식으로 이해하더라도}\srcfn{9)}{실제로 쓰는 것은 정수적 원소들이 통상적인 계산법칙 아래 덧셈, 뺄셈, 곱셈으로 닫혀서 환을 이룬다는 것, 그리고 이 환에서 곱이 사라지면 한 인자가 사라져서 원소쌍을 첨가하는 몫구성으로 체까지 확대할 수 있다는 것뿐이다. 반면 모듈의 기저표현은 쓰지 않는다. 따라서 예를 들어 계수들이 모든 정수적 대수적 수의 정역에 속할 때에도 (7)이 성립한다. (E. N.)}, \emph{몫표현
\begin{equation}
\Gmod=\Amod/(d)
\tag{7}
\end{equation}
이 여전히 성립한다. 여기서 \(d\)는 \(\Amod\)에서 얻는 항등적으로 영이 아닌 임의의 \(\varrho\)차 행렬식이다.}

이 확대 아래에서도 랭크, 기초모듈, 모듈몫의 개념은 유지된다. 다만 이제 모듈몫은 주아이디얼일 필요가 없다. 또한 기저표현을 제외하면 정리 I도 유지된다.

이제
\[
\begin{aligned}
a_1(\xi)&=a_{11}\xi_1+\cdots+a_{1k}\xi_k,\ldots,\\
a_\varrho(\xi)&=a_{\varrho1}\xi_1+\cdots+a_{\varrho k}\xi_k
\end{aligned}
\]
가 \(\Amod\)에 속하는 \(\varrho\)개의 선형독립인 선형형식이라고 하자. 이들은 일반적으로 모듈기저를 이루지 않는다. 부정원 \(\xi\)들의 번호를
\[
d=|a_{ij}|\ne0,\qquad i,j=1,2,\ldots,\varrho
\]
가 되도록 정하자. 그러면 다음 특별한 모양의 \(\varrho\)개 선형형식이 \(\Amod\)에 속한다.
\begin{equation}
d\xi_\lambda+b_{\lambda,\varrho+1}\xi_{\varrho+1}
+\cdots+b_{\lambda,k}\xi_k\zero{\Amod},
\qquad(\lambda=1,2,\ldots,\varrho).
\tag{8}
\end{equation}
그러나 \(\Amod\)은 \(\xi_{\varrho+1},\ldots,\xi_k\)에만 의존하는 선형형식
\[
c_{\varrho+1}\xi_{\varrho+1}+\cdots+c_k\xi_k
\]
을 포함할 수 없다. 그렇다면 적어도 하나의 \((\varrho+1)\)차 행렬식 \(d\,c_\alpha\)가 영이 아닐 것이기 때문이다.

이제 몫 \(\Amod/(d)\)는 \(\Amod\)의 같은 랭크를 가진 인자이므로 \(\Gmod\)로 나누어진다. 그러나 역도 성립한다. 각 \(g(\xi)\zero{\Gmod}\)에 대하여 정의상
\[
c\,g(\xi)\zero{\Amod},\qquad c\ne0,
\]
따라서
\[
dc\,g(\xi)\zero{\Amod}
\]
이다. 한편 (8)에 따라
\[
d\,g(\xi)=g_{\varrho+1}\xi_{\varrho+1}+\cdots+g_k\xi_k
\zero{\Amod}.
\]
그러므로
\[
cg_{\varrho+1}\xi_{\varrho+1}+\cdots+cg_k\xi_k\zero{\Amod}.
\]
방금 지적한 사실에 따라 \(cg_{\varrho+1}=0,\ldots,cg_k=0\)이고, \(c\ne0\)이므로 \(g_{\varrho+1}=0,\ldots,g_k=0\)이다. 따라서 \(d\,g(\xi)\zero{\Amod}\)이며 (7)이 증명되었다.

\(d\,g(\xi)\zero{\Amod}\)은 두 모듈
\[
(a_1(\xi),\ldots,a_\varrho(\xi)),\qquad
(g(\xi),a_1(\xi),\ldots,a_\varrho(\xi))
\]
이 같은 랭크 \(\varrho\)를 갖는다는 사실에서도 직접 따라온다. 실제로 \(g(\xi)=c_1\xi_1+\cdots+c_k\xi_k\)로 두면 \((\varrho+1)\)차 행렬식
\[
\begin{vmatrix}
a_1(\xi)&a_{11}&\cdots&a_{1\varrho}\\
\vdots&\vdots&&\vdots\\
a_\varrho(\xi)&a_{\varrho1}&\cdots&a_{\varrho\varrho}\\
g(\xi)&c_1&\cdots&c_\varrho
\end{vmatrix}
\]
이 사라진다. 그러나 위에서 준 증명은 제4절에서 식 (30)에 이어 다시 나타난다.

\section*{\S{}2. 노름과 모듈의 분해정리}

이제 다시 유리정수 또는 \(P\)를 계수체로 갖는 한 부정원의 다항식에 관한 선형형식 모듈을 다룬다.

\medskip
\noindent\textbf{정의 III.} \emph{통상처럼 \(\xi\)의 선형형식 가운데 \(\Amod\)을 법으로 서로 합동인 것들을 한 잉여류로 묶자. 기호 \(\Gmod\mid\Amod\)은 \(\Amod\)에 관한 \(\Gmod\)의 잉여류계, 즉 \(\Gmod\)의 원소들로 이루어진 \(\Amod\)의 모든 잉여류의 계를 뜻한다. \(\Amod\)에 대한 \(\Gmod\)의 노름, 기호로 \(\Norm(\Gmod\mid\Amod)\), 은 \(\Gmod\)에서 \(\Amod\)으로 넘어가는 치환의 행렬식으로 정의한다. 즉 \(\Amod\)의 임의의 선형독립인 모듈기저를 \(\Amod\)의 기초모듈 \(\Gmod\)의 그러한 기저로 나타내는 치환의 행렬식이다.}

(1), (3), 또는 정리 II 뒤의 주에서 다음을 얻는다.
\begin{equation}
\Norm(\Gmod\mid\Amod)=e_1e_2\cdots e_\varrho=d_\varrho,
\qquad
\Gmod=\Amod/\bigl(\Norm(\Gmod\mid\Amod)\bigr).
\tag{9}
\end{equation}

(1), (3), (9)에서 잘 알려진 방식으로 다음이 따라온다. 유리정수의 경우 \(\Norm(\Gmod\mid\Amod)\)은 \(\Amod\)에 관한 \(\Gmod\)의 잉여류 수를 나타낸다. 한 부정원의 다항식인 경우에는 잉여류의 수가 일반적으로 무한이지만, \(\Norm(\Gmod\mid\Amod)\)의 차수는 \emph{\(P\)에 관하여 선형독립인 잉여류의 수}와 같다. \(\Bmod\)이 \(\Amod\)의 같은 랭크를 가진 인자이면, \(\Bmod\)은 그러한 한 잉여류의 대표원과 함께 그 잉여류의 모든 원소를 포함한다. 따라서 \(\Bmod\)은 \(\Amod\)에 유한개의 잉여류 또는 그 선형결합들을 덧붙여 얻는다. 여기서 곧바로 다음이 따라온다.

\medskip
\noindent\textbf{정리 IV.} \emph{\(\Bmod\)이 \(\Amod\)의 같은 랭크를 가진 인자라서 두 기초모듈이 같고, 또한
\[
\Norm(\Gmod\mid\Bmod)=\Norm(\Gmod\mid\Amod)
\]
이면 \(\Bmod=\Amod\)이다.}

노름의 분해와 모듈의 분해 사이의 관계에 대해서는 다음이 성립한다.

\medskip
\noindent\textbf{정리 V.} \emph{
\begin{equation}
\Norm(\Gmod\mid\Amod)=rs,\qquad(r,s)=1
\tag{10}
\end{equation}
이라고 하자. 그러면 \(\Amod\)의 같은 랭크를 가진 인자 \(\Rmod\)과 \(\Smod\)이 각각 유일하게 존재하여
\[
r=\Norm(\Gmod\mid\Rmod),\qquad
s=\Norm(\Gmod\mid\Smod)
\]
가 된다. \(\Rmod,\Smod\)은
\[
\Rmod=\Amod/(s),\qquad
\Smod=\Amod/(r)
\]
로 정의되며, \(\Amod\)은 \(\Rmod,\Smod\)의 최소공배모듈과 같다.}\srcfn{10)}{최소공배모듈 \([\Rmod,\Smod]\)은 \(\Rmod\)로도 \(\Smod\)으로도 나누어지는 모든 선형형식의 전체이다. 이에 대응하여 최대공약모듈 \((\Rmod,\Smod)\)은 \(\Rmod\)의 한 원소와 \(\Smod\)의 한 원소의 합으로 표현되는 모든 선형형식의 전체이다.}
\emph{\(r_\varrho=(r,e_\varrho)\), \(s_\varrho=(s,e_\varrho)\)로 두면 다음 역관계가 성립한다.
\begin{equation}
(s_\varrho)=\Amod/\Rmod,\quad
\Rmod=\Amod/(s_\varrho);\qquad
(r_\varrho)=\Amod/\Smod,\quad
\Smod=\Amod/(r_\varrho).
\tag{11}
\end{equation}}

실제로 일반적으로 \(r_i=(r,e_i)\), \(s_i=(s,e_i)\)로 두면 (9), (10)에 따라
\[
(r_i,s_i)=1,\qquad e_i=r_is_i,\qquad
r=r_1\cdots r_\varrho,\qquad s=s_1\cdots s_\varrho,
\]
이고 각 \(r_i\), 또는 \(s_i\)는 다음 \(r_{i+1}\), 또는 \(s_{i+1}\)을 나눈다. 따라서
\[
\Rmod=(r_1\eta_1,\ldots,r_\varrho\eta_\varrho),\qquad
\Smod=(s_1\eta_1,\ldots,s_\varrho\eta_\varrho)
\]
로 두면 \(\Rmod,\Smod\)은 \(\Amod\)의 같은 랭크를 가진 인자가 된다. \(r_i,s_i\)가 서로소이므로 \(\Amod\)은 \(\Rmod,\Smod\)의 최소공배모듈이고, 또한
\[
\Norm(\Gmod\mid\Rmod)=r_1\cdots r_\varrho=r,\qquad
\Norm(\Gmod\mid\Smod)=s_1\cdots s_\varrho=s
\]
이다.

더 나아가
\[
s_\varrho\Rmod\zero{\Amod},\qquad
r_\varrho\Smod\zero{\Amod}.
\]
\(c\Rmod\zero{\Amod}\)이면 \(cr_\varrho\eta_\varrho\zero{\Amod}\), 따라서 \(c\zero{(s_\varrho)}\)이다. 이로써 \((s_\varrho)=\Amod/\Rmod\)이 증명되고, 마찬가지로 \((r_\varrho)=\Amod/\Smod\)이 증명된다. 또
\[
b(\eta)=b_1\eta_1+\cdots+b_\varrho\eta_\varrho
\zero{\Amod/(s_\varrho)}
\]
이면 모든 \(r_i\)가 \(s_\varrho\)와 서로소이므로 \(b_i\zero{(r_i)}\)이다. 따라서 \(\Rmod=\Amod/(s_\varrho)\)이고, 마찬가지로 \(\Smod=\Amod/(r_\varrho)\)이다. 이로써 역관계 (11)이 증명되었다. 특수한 경우 \(r=1\)이면 (11)은 (5)가 된다.

이제 \(\Rmod,\Smod\)이 \emph{유일하게 정해짐}을 보여야 한다. \(\overline{\Rmod},\overline{\Smod}\)이 정리 V의 조건을 만족하는 모듈들이고, \(\bar r_i,\bar s_i\)가 각각의 초등인자라고 하자. \(\bar r_i,\bar s_i\)는 \(e_i\)의 인자이고
\[
r=\bar r_1\cdots\bar r_\varrho,\qquad
s=\bar s_1\cdots\bar s_\varrho,\qquad
(\bar r_i,\bar s_i)=1
\]
이므로
\[
e_i=\bar r_i\bar s_i,\qquad
\bar r_i=(r,e_i)=r_i,\qquad
\bar s_i=(s,e_i)=s_i
\]
이다. 따라서 \(\overline{\Rmod},\overline{\Smod}\)은 \(\Rmod,\Smod\)과 초등인자와 기초모듈이 각각 같다. 이제 단모듈러치환 \(\eta=Q(\xi)\)로 새 부정원 \(\xi\)들을 도입하면
\[
\overline{\Rmod}=(r_1\xi_1,\ldots,r_\varrho\xi_\varrho)
\]
이고
\[
\Amod=\bigl(
r_1s_1(q_{11}\xi_1+\cdots+q_{1\varrho}\xi_\varrho),\ldots,
r_\varrho s_\varrho(q_{\varrho1}\xi_1+\cdots+q_{\varrho\varrho}\xi_\varrho)
\bigr).
\]
\(\Amod\zero{\Rmod}\)에서
\[
r_is_i(q_{i1}\xi_1+\cdots+q_{i\varrho}\xi_\varrho)\zero{\Rmod}
\]
이고, 따라서 \(r_is_iq_{ij}\zero{(r_i)}\)이다. \((r,s)=1\)이므로 \(r_iq_{ij}\zero{(r_i)}\), 즉 \(\overline{\Rmod}\zero{\Rmod}\)이다. 그런데
\[
\Norm(\Gmod\mid\Rmod)=\Norm(\Gmod\mid\overline{\Rmod})
\]
이므로 정리 IV에 따라 \(\Rmod=\overline{\Rmod}\)이다. 마찬가지로 \(\Smod=\overline{\Smod}\)이다. 이로써 정리 V가 모두 증명되었다.

\section*{\S{}3. 다항식의 아이디얼}

\(\overline{\mideal}\)이 추상적으로 정의된 임의의 체 \(P\)를 계수체로 갖는 \(n\)개 부정원 \(y_1,\ldots,y_n\)의 \emph{다항식 아이디얼}이라고 하자. 즉 \(\overline{\mideal}\)은 \(\Phi_1(y),\Phi_2(y)\)와 함께 차 \(\Phi_1(y)-\Phi_2(y)\)를 포함하고, \(\Phi(y)\)와 함께 \(P\)의 계수를 갖는 임의의 다항식 \(A(y)\)에 대하여 \(A(y)\Phi(y)\)도 포함한다\srcfn{11)}{개념상 자연스러운 \glqq{}아이디얼\grqq{}이라는 명칭은 예전에 널리 쓰던 \glqq{}모듈\grqq{} 또는 \glqq{}형식모듈\grqq{} 대신 쓰인다. 여기서는 이미 선형형식 모듈과 구별하기 위해서도 이 명칭이 필요하다. (E. N.)}.

부정원들을 계수 자체가 부정원인 변환
\begin{equation}
y=U(x),\qquad\text{예를 들어}\qquad
\begin{cases}
y_1=x_1,\\
y_2=u_{21}x_1+x_2,\\
\cdots\\
y_n=u_{n1}x_1+\cdots+u_{n,n-1}x_{n-1}+x_n
\end{cases}
\tag{12}
\end{equation}
에 놓고, 부정원 \(u_{\mu\nu}\)들을 다항식의 계수영역에 첨가하자. 그러면 계수영역은 체 \(P(u)\)로 바뀐다. 이 첨가로 \(\overline{\mideal}\)이 \(\overline{\mideal}_{(u)}\)로 바뀐다고 하자. 그러면 \(\overline{\mideal}_{(u)}\)은 \(\overline{\mideal}\)의 다항식 \(\Phi_\lambda(y)\)들과 함께, 매번 유한개의 \(\Phi_\lambda\)를 취한 모든 선형결합
\[
\sum_\lambda\alpha_\lambda(u)\Phi_\lambda(y)
\]
도 포함한다. 여기서 \(\alpha_\lambda(u)\)는 \(P\)의 계수를 갖는 \(u_{\mu\nu}\)들의 유리함수이다. 알맞은 \(u\)의 다항식을 곱하면 일반성을 잃지 않고 \(\alpha_\lambda(u)\)들을 \(u\)의 거듭제곱곱이라고 할 수 있다. 따라서
\begin{equation}
\sum_\lambda\alpha_\lambda(u)\Phi_\lambda(y)
\zero{\overline{\mideal}_{(u)}},
\qquad
\Phi_\lambda(y)\zero{\overline{\mideal}},
\quad\text{그리고 그 역도 성립한다.}
\tag{13}
\end{equation}
(12)에 따라 \(\overline{\mideal}_{(u)}\)은 \(P(u)\)를 계수체로 갖는 \(x_1,\ldots,x_n\)의 다항식 아이디얼 \(\mideal\)로 옮겨진다.
\begin{equation}
\overline{\mideal}_{(u)}=\mideal
\qquad\text{단, }y=U(x).
\tag{14}
\end{equation}
관계 (14), (13)을 결합하면 다음을 얻는다.
\begin{equation}
\begin{gathered}
F(x)\zero{\mideal},\qquad
F(x)=\Phi(y)=\sum_i\alpha_i(u)\Phi_i(y)
=\sum_i\alpha_i(u)F_i(x)\\
\Longrightarrow\qquad F_i(x)\zero{\overline{\mideal}}.
\end{gathered}
\tag{15}
\end{equation}
반대로 (14), (15)에서 다시 (13)을 얻는다.

\medskip
\noindent\textbf{정의 IV.} \emph{\(P(u)\)를 계수체로 갖는 \(x_1,\ldots,x_n\)의 다항식 아이디얼 \(\mideal\) 가운데 관계 (15)를 만족하는 것, 즉 \(y=U(x)\)에 따라 \(\overline{\mideal}_{(u)}\)에서 생기는 것을 변환된 아이디얼이라 한다.}

변환된 아이디얼의 특징은 정칙 다항식의 존재이다. \emph{\(r\)차 다항식이 영이 아닌 계수의 항 \(x_i^r\)을 포함하면 \(x_i\)에 관하여 정칙이라고 한다.} 다항식 \(F_i(x)\)들이 \(x_1\)에 관하여 정칙임을 알 수 있다. 이하에서는 주로 이 변환된 아이디얼들을 \(x\)의 특정한 거듭제곱곱들에 관한 선형형식 모듈로 볼 것이다. 이를 위해 나중에 기초모듈로 옮겨갈 기초아이디얼의 개념을 정해야 한다. 먼저 이제부터 계속 쓸 표기를 정한다.

\medskip
\noindent\textbf{표기.} \emph{\(F^{(i)},f^{(i)},a^{(i)},\ldots\)는 언제나 \(x_1,\ldots,x_{i-1}\)과 무관한 \(x\)의 다항식을 뜻한다.}

\medskip
\noindent\textbf{정의 V.} \emph{각 아이디얼 \(\mideal\)에는 \(n\)개의 기초아이디얼 \(\gideal_0,\gideal_1,\ldots,\gideal_{n-1}\)이 다음처럼 정의된다. \((i-1)\)단계 기초아이디얼 \(\gideal_{i-1}\)은, 일반적으로 \(G(x)\)에 따라 달라지는 다항식 \(b^{(i)}\)가 존재하여
\begin{equation}
b^{(i)}G(x)\zero{\mideal},\qquad b^{(i)}\ne0
\tag{16}
\end{equation}
이 되는 다항식 \(G(x)\)만을, 그리고 그 모두를 포함한다.}

\(\gideal\)들이 실제로 아이디얼이라는 증명은 식 (4) 뒤에서 \(\Gmod\)의 모듈 성질을 증명한 것과 정확히 대응하며, (16)은 (4)의 유사물이다. 각 다항식 \(b^{(i+1)}\)은 동시에 \(b^{(i)}\)이므로 \(\gideal_i\)는 \(\gideal_{i-1}\)로 나누어진다.

기초아이디얼에 대해서는 다음이 성립한다.

\medskip
\noindent\textbf{정리 VI.} \emph{변환된 아이디얼의 기초아이디얼들도 변환된 아이디얼이다}\srcfn{12)}{정리 VI은 제7절에서 처음 쓴다.}.

증명은 잘 알려진 데데킨트--메르텐스 정리\srcfn{13)}{Dedekind, \emph{Über einen arithmetischen Satz von Gauß}, \emph{Mitt. d. deutsch. math. Ges. zu Prag} (1892); F. Mertens, \emph{Über einen algebraischen Satz}, \emph{Ber. d. Ak. d. Wissensch. Wien} 101 (1892), 1560--1566쪽. 크로네커가 가우스 정리를 부정계수 다항식으로 확장한 것은 이 정리의 직접적인 귀결이다.}에 바탕을 둔다. \(a,b\)가 부정원이고
\[
\begin{aligned}
&\sum a_{i_1,\ldots,i_s}z_1^{i_1}\cdots z_s^{i_s}\,
\sum b_{j_1,\ldots,j_s}z_1^{j_1}\cdots z_s^{j_s}\\
&\hspace{8em}=
\sum c_{k_1,\ldots,k_s}z_1^{k_1}\cdots z_s^{k_s},
\end{aligned}
\]
\[
\sum i\le\nu,\qquad
\sum j\le\mu,\qquad
\sum k\le\nu+\mu
\]
라고 하자. 또 \(\Amod,\Bmod,\mathfrak C\)가 각각 \(a,b,c\)들로부터 유리정수에 의해 생성되는 모듈이라고 하자. 그러면 어떤 지수 \(q\)가 존재하여 항등식
\begin{equation}
\Amod^q\Bmod=\Amod^{q-1}\mathfrak C
\tag{17}
\end{equation}
이 성립한다. 여기서 \(\Amod^q\)는 \(a\)들의 \(q\)차 거듭제곱곱들과 그 정수계수 선형결합들로 이루어진다.

정리 VI을 증명하기 위해
\begin{equation}
\begin{cases}
G(x)=\Gamma(y)=\sum_\lambda\alpha_\lambda(u)\Gamma_\lambda(y)
=\sum_\lambda\alpha_\lambda(u)G_\lambda(x),\\
b^{(i)}(x)=\varphi(y)=\sum_\mu\beta_\mu(u)\varphi_\mu(y)
\end{cases}
\tag{18}
\end{equation}
로 두자. 여기서 \(\alpha_\lambda(u),\beta_\mu(u)\)는 \(u\)의 거듭제곱곱들이다. 그러면 (16), (14), (18)에 따라
\begin{equation}
\sum_\mu\beta_\mu(u)\varphi_\mu(y)\,
\sum_\lambda\alpha_\lambda(u)\Gamma_\lambda(y)
\zero{\overline{\mideal}_{(u)}},
\qquad
\sum_\mu\beta_\mu(u)\varphi_\mu(y)\ne0.
\tag{19}
\end{equation}
왼쪽은 \(u\)에 관한 두 다항식의 곱이고, 그 계수는 다항식 \(\varphi_\mu(y),\Gamma_\lambda(y)\)이다. (19)의 오른쪽에 따라 곱다항식의 \(u\)-계수들은 \(\overline{\mideal}\)에 속하는 다항식이다. 따라서 (17)에서 \(a,b,c\)를 이 \(y\)의 다항식들로 바꾸면
\[
\left\{\sum_\mu\beta_\mu(u)\varphi_\mu(y)\right\}^{q}
\Gamma_\lambda(y)\zero{\overline{\mideal}_{(u)}
}
\]
를 얻는다. (18)에 따라 이는
\begin{equation}
b^{(i)}(x)^qG_\lambda(x)\zero{\mideal},
\qquad
G_\lambda(x)\zero{\gideal_{i-1}}
\tag{20}
\end{equation}
과 같다. 관계 (16), (18), (20)은 \emph{기초아이디얼 \(\gideal_{i-1}\)들이 실제로 변환된 아이디얼}임을 보여 준다.

\section*{\S{}4. 다항식 아이디얼과 선형형식 모듈의 관계}

언제나 변환된 것이라고 가정하는 다항식 아이디얼 \(\mideal\)을 선형형식 모듈 \(\Mmod_{i-1}\), \(i=1,\ldots,n\), 로 보기 위해 각 다항식 \(F(x)\)를 \(x_1,\ldots,x_{i-1}\)의 거듭제곱곱 \(\xi_\lambda\)들에 관한 선형형식으로 본다.
\[
F(x)=\sum_\lambda a_\lambda^{(i)}\xi_\lambda.
\]
제3절의 표기에 따라 \(a_\lambda^{(i)}\)는 \(x_i,\ldots,x_n\)의 다항식이다. 아이디얼의 정의에서 \(\Mmod_{i-1}\)이 이 정수적 원소 \(a^{(i)},b^{(i)},\ldots\)에 관하여 모듈 성질을 가짐이 곧바로 따라온다. 또한 \(x_1,\ldots,x_{i-1}\)의 무한히 많은 거듭제곱곱 \(\xi\)는 오직 \emph{선형}으로 나타나므로, 서로 선형독립이라는 사실 때문에 \(\Mmod_{i-1}\)에 대하여 부정원의 역할을 한다. 각 모듈 \(\Mmod_{i-1}\)은 무한히 많은 부정원 \(\xi\)를 포함하지만, 한 선형형식에는 그 가운데 유한개만 나타난다. 마찬가지로 각 기초아이디얼 \(\gideal_{i-1}\)을 선형형식 모듈 \(\Gmod_{i-1}\)로 보며, 부정원은 다시 \(x_1,\ldots,x_{i-1}\)의 거듭제곱곱이다. \(i=1\)일 때에는 단위원에 대응하는 부정원 \(\xi_0\) 하나만 있다.

이하의 모든 내용은 다음 정리에 따라, 이처럼 부정원 \(\xi\)가 무한히 많아도 유한개의 부정원에 관한 선형형식 모듈로 제한할 수 있다는 사실에 바탕을 둔다.

\medskip
\noindent\textbf{정리 VII.} \emph{잉여류계 \(\Gmod_{i-1}\mid\Mmod_{i-1}\)은 잉여류계 \(\Gmod^*_{i-1}\mid\Mmod^*_{i-1}\)과 동형이다. 여기서 \(\Mmod^*_{i-1}\)은 유한개의 부정원에 관한 모듈이고 \(\Gmod^*_{i-1}\)은 그 기초모듈이다. \(x_{i+1},\ldots,x_n\)을 \(P(u)\)에 첨가하면 \(\Mmod_{i-1}\)은 \(1\)과 다른 초등인자를 유한개만 가지며, 이들은 \(\Mmod^*_{i-1}\)의 \(1\)과 다른 초등인자들이다.}

여기서 \(\Mmod_{i-1}\)의 초등인자는 주 8과 같이 정의한다. 정리 VII의 동형은 다음 정의에 따른다.

\medskip
\noindent\textbf{정의 V.} \emph{두 잉여류계를 일대일로 대응시킬 수 있고, 두 류의 차에는 대응하는 두 류의 차가 대응하며, 한 류와 한 정수적 원소의 곱에는 대응하는 류와 같은 정수적 원소의 곱이 대응하면 두 잉여류계를 동형이라고 한다.}

정리 VII은 완전귀납법으로 증명한다. \(i=1\)일 때에는 분명하다. \(\Mmod_0\)은 \(x\)의 0차 거듭제곱곱, 즉 단위원에 대응하는 부정원 \(\xi_0\) 하나에 관한 선형형식 모듈이고, 정수적 원소들은 \(x_1,\ldots,x_n\)의 모든 다항식이다. 기초아이디얼 \(\gideal_0\)은 \(x_1,\ldots,x_n\)의 모든 다항식으로 이루어진 단위원아이디얼이고, \(\Gmod_0\)은 \(\Mmod_0\)의 기초모듈인 모듈 \((\xi_0)\)과 같다. 따라서 \(\Gmod^*_0\mid\Mmod^*_0\)은 \(\Gmod_0\mid\Mmod_0\)과 같다. 끝으로 랭크 \(1\)인 \(\Mmod_0\)은 \(1\)과 다른 초등인자를 기껏해야 하나만 갖는다.

먼저 \(i=1\)에서 \(i=2\)로 넘어간다. 이때 \(\Gmod_1\mid\Mmod_1\)의 구조를 파악한 뒤 일반적인 귀납단계를 수행할 것이다. 우선 \(\Gmod_0\mid\Mmod_0\)에 대하여 \(x_1\)의 차수가 제한된 대표원계를 택할 수 있음을 보인다. 그러면 \(\gideal_1\)이 \(\gideal_0\)으로 나누어진다는 사실에 따라 \(\Gmod^*_1\)의 부정원은 유한개가 된다.

제3절에 따라 변환된 아이디얼 \(\mideal\)은 \(x_1\)에 관하여 정칙인 다항식 \(C^{(1)}(x)\)을 적어도 하나 포함한다. 그 차수를 \(k\)라 하자. 따라서 \(\Mmod_0\)은 선형형식 \(C^{(1)}\xi_0\)을 포함한다. \(C^{(1)}(x)\)의 정칙성 때문에 모든 다항식 \(H(x)\)은
\begin{equation}
H(x)\equiv b_0^{(2)}+b_1^{(2)}x_1+\cdots+
b_{k-1}^{(2)}x_1^{k-1}
\pmod{C^{(1)}(x)}
\tag{21}
\end{equation}
로 표현된다. 그러므로 \(C^{(1)}\xi_0\)이 \(\Mmod_0\)에 속한다는 사실에 따라, \(\Gmod_0\mid\Mmod_0\)의 대표원계를 \(x_1\)에 관한 차수가 \(k\)에 미치지 않는 것들로 실제로 택할 수 있다.

특히 \(\gideal_1\)의 다항식 \(G(x)\)와 \(\mideal\)의 다항식 \(F(x)\)에 대하여
\begin{equation}
\left\{
\begin{aligned}
G(x)&\equiv g_0^{(2)}+g_1^{(2)}x_1+\cdots+g_{k-1}^{(2)}x_1^{k-1}
&&\pmod{C^{(1)}(x)},\\
F(x)&\equiv a_0^{(2)}+a_1^{(2)}x_1+\cdots+a_{k-1}^{(2)}x_1^{k-1}
&&\pmod{C^{(1)}(x)}
\end{aligned}\right.
\tag{22}
\end{equation}
로 두자. \(\Gmod^*_1\)은 모든 선형형식
\[
g(\xi)=g_0^{(2)}\xi_0+\cdots+g_{k-1}^{(2)}\xi_{k-1}
\]
의 계, \(\Mmod^*_1\)은 이에 대응하는 모든
\[
a(\xi)=a_0^{(2)}\xi_0+\cdots+a_{k-1}^{(2)}\xi_{k-1}
\]
의 계라고 하자. \(\gideal_1,\mideal\)의 아이디얼 성질로부터 \(\Gmod^*_1,\Mmod^*_1\)은 정수적 원소 \(b^{(2)},c^{(2)},\ldots\)에 관한, \(\xi_0,\ldots,\xi_{k-1}\)의 선형형식 모듈임이 따라온다. \(C^{(1)}(x)\)가 생성하는 주아이디얼도 이 정수적 원소들에 관한 모듈
\[
\mathfrak C_1=(\zeta_0,\zeta_1,\ldots,\zeta_\nu,\ldots)
\]
을 주며, 여기서 \(\zeta_\nu=x_1^\nu C^{(1)}(x)\)이다. \(\zeta_\nu\)들은 정수적 원소 \(b^{(2)},c^{(2)},\ldots\)에 관하여 서로 선형독립이고, 유한개의 \(\xi\)들과도 선형독립이다. \(\mathfrak C_1\)이 \(\Mmod_1\), 따라서 \(\Gmod_1\)로 나누어진다는 사실에서 \(\Gmod^*_1\)도 \(\Gmod_1\)로, \(\Mmod^*_1\)도 \(\Mmod_1\)로 각각 나누어짐이 따라온다. (22)를 함께 고려하면
\begin{equation}
\Gmod_1=(\Gmod^*_1,\mathfrak C_1),\qquad
\Mmod_1=(\Mmod^*_1,\mathfrak C_1)
\tag{23}
\end{equation}
을 얻는다.

(23)과 \(\xi,\zeta\)의 선형독립성에서 \(i=2\)에 대한 정리 VII이 다음처럼 따라온다. 먼저 \(\Gmod_1\)이 \(\Mmod_1\)로 나누어지므로
\[
\Gmod_1\mid\Mmod_1=\Gmod^*_1\mid\Mmod_1.
\]
\(\Gmod^*_1\mid\Mmod^*_1\)과의 동형을 보이기 위해, \(\Gmod^*_1\)의 어떤 선형형식 \(g^*(\xi)\)들이 \(\Gmod^*_1\mid\Mmod^*_1\)의 대표원계를 이룬다고 하자. \(\xi,\zeta\)의 선형독립성 때문에 \(g^*(\xi)\equiv0\pmod{\Mmod_1}\)이면 \(g^*(\xi)\equiv0\pmod{\Mmod^*_1}\)이다. 따라서 \(g^*(\xi)\)들은 \(\Mmod_1\)을 법으로도 서로 합동이 아니고, \(\Gmod_1\mid\Mmod_1\)의 대표원계를 이룬다. 이 대표원계를 통해 \(\Gmod^*_1\mid\Mmod_1\)과 \(\Gmod^*_1\mid\Mmod^*_1\) 사이에 생기는 대응은 정의 V의 의미에서 동형이다.

똑같이 \(b^{(2)}g(\xi)\equiv0\pmod{\Mmod_1}\), \(b^{(2)}\ne0\)이면 \(b^{(2)}g(\xi)\equiv0\pmod{\Mmod^*_1}\), \(b^{(2)}\ne0\)이다. 이 조건은 \(\Gmod^*_1\)에 속하는 모든 \(g(\xi)\)에 대하여, 그리고 오직 그 형식들에 대해서만 성립한다. (23)에 따라 \(\Gmod^*_1\)은 기초아이디얼 \(\gideal_1\)의 다항식 가운데 동시에 \(\xi\)의 선형형식인 것들 전체이기 때문이다. 따라서 \(\Gmod^*_1\)은 \(\Mmod^*_1\)의 기초모듈로 특징지어진다.

끝으로 \(x_3,\ldots,x_n\)을 \(P(u)\)에 첨가하면
\begin{equation}
\left\{
\begin{aligned}
\Gmod^*_1&=(\eta_1,\ldots,\eta_\rho),&
\Mmod^*_1&=(e_1\eta_1,\ldots,e_\rho\eta_\rho),\\
\Gmod_1&=(\eta_\rho,\ldots,\eta_1,\zeta_0,\ldots,\zeta_\nu,\ldots),&
\Mmod_1&=(e_\rho\eta_\rho,\ldots,e_1\eta_1,\zeta_0,\ldots,\zeta_\nu,\ldots)
\end{aligned}\right.
\tag{24}
\end{equation}
이고, 따라서
\[
(e_\rho)=\Mmod_1/\Gmod_1
=\Mmod_1/(\Mmod_1,\eta_\rho),\qquad
(e_{\rho-1})=(\Mmod_1,\eta_\rho)/\Gmod_1,\quad\text{등등}
\]
이다. 주 8까지 고려하면 \(\Mmod_1\)의 초등인자는
\[
e_\rho,e_{\rho-1},\ldots,e_1,1,\ldots,1,\ldots
\]
임을 알 수 있다. 이로써 \(i=2\)에 대한 정리 VII의 모든 부분이 증명되었다.

(22), (23)은 \(C^{(1)}(x)\)가 \(x_1\)에 관하여 정칙이라는 사실만을 썼다는 점을 주목하자. 따라서 (12) 대신 특수한 변환
\begin{equation}
y_1=x_1,\qquad
y_2=u_{21}x_1+z_2,\quad\ldots,\quad
y_n=u_{n1}x_1+z_n
\tag{25}
\end{equation}
을 기초로 해도 이 식들과 정리 VII로 이어지는 논증은 그대로 성립한다. 이를
\begin{equation}
\left\{
\begin{array}{l}
z=U_1(x),\quad\text{또는}\\[1mm]
z_2=x_2,\quad z_3=u_{32}x_2+x_3,\quad\ldots,\quad
z_n=u_{n2}x_2+\cdots+u_{n,n-1}x_{n-1}+x_n
\end{array}\right.
\tag{26}
\end{equation}
과 합성하면 다시 (12)를 얻는다. \(\widetilde{\mideal}_{(u)}\)가 (25)에 따라 \(\widetilde{\mideal}\)로 옮겨지고, \(\widetilde{\gideal}_1,\widetilde{\Gmod}_1,\ldots\)가 \(\widetilde{\mideal}\)에 대하여 \(\gideal_1,\Gmod_1,\ldots\)과 같은 뜻을 갖는다면
\begin{equation}
\widetilde{\Gmod}_1=(\widetilde{\Gmod}^{*}_1,\widetilde{\mathfrak C}_1),
\qquad
\widetilde{\Mmod}_1=(\widetilde{\Mmod}^{*}_1,\widetilde{\mathfrak C}_1)
\tag{27}
\end{equation}
이다.

(27)은 (26)의 역변환에 따라 (23)에서 단순히 생긴다. 실제로 이는 \(\mideal,C^{(1)}(x)\), 따라서 \(\mathfrak C_1,\Mmod_1\)에 대해서 분명히 성립한다. \(\gideal_1\)에 대해서도 성립한다. 관계
\[
b^{(2)}(x)G(x)\equiv0\pmod{\mideal},\quad b^{(2)}\ne0,
\]
\[
\widetilde b^{(2)}(z)\widetilde G(x,z)
\equiv0\pmod{\widetilde\mideal},\quad
\widetilde b^{(2)}\ne0
\]
은 \(z=U_1(x)\)에 따라 서로를 함의하기 때문이다. 그러므로 (26)에 따라 \(\widetilde\gideal_1=\gideal_1\), 따라서 \(\widetilde\Gmod_1=\Gmod_1\)이다. (22)가 주는 \(\Gmod_1^*,\Mmod_1^*\)의 정의에 따라 \(\widetilde\Gmod_1^*,\widetilde\Mmod_1^*\), 따라서 표현 (27) 전체에도 같은 사실이 성립한다.

일반적인 귀납단계를 위해 다음을 가정하자.
\begin{equation}
\left\{
\begin{aligned}
\Gmod_{i-1}&=(\Gmod_{i-1}^*,\mathfrak C_{i-1}),&
\Mmod_{i-1}&=(\Mmod_{i-1}^*,\mathfrak C_{i-1}),\\
\mathfrak C_{i-1}&=(\zeta_0,\zeta_1,\ldots,\zeta_r,\ldots).
\end{aligned}\right.
\tag{28}
\end{equation}
여기서 \(\Gmod^*_{i-1},\Mmod^*_{i-1}\)은 정수적 원소 \(b^{(i)},c^{(i)},\ldots\)에 관한, 유한개의 부정원 \(\xi\)의 선형형식 모듈이다. 이 \(\xi\)들은 \(x_1,\ldots,x_{i-1}\)의 특정한 거듭제곱곱들이다. 또한 \(\zeta\)들은 이 정수적 원소들에 관하여 서로 선형독립이고 \(\xi\)들과도 선형독립이라고 하자. (12) 대신 특수한 변환
\begin{equation}
\begin{aligned}
y_1&=x_1,\quad\ldots,\quad
y_{i-1}=u_{i-1,1}x_1+\cdots+x_{i-1},\\
y_i&=u_{i,1}x_1+\cdots+u_{i,i-1}x_{i-1}+z_i,\quad\ldots,\quad
y_n=u_{n,1}x_1+\cdots+u_{n,i-1}x_{i-1}+z_n
\end{aligned}
\tag{29}
\end{equation}
을 기초로 할 때에도 (28)에 대응하는 표현과 \(\xi,\zeta\)의 선형독립성이 성립한다고 하자. 이 변환은
\[
z=U_{i-1}(x),
\]
즉
\[
z_i=x_i,\quad z_{i+1}=u_{i+1,i}x_i+x_{i+1},\quad\ldots,\quad
z_n=u_{n,i}x_i+\cdots+u_{n,n-1}x_{n-1}+x_n
\]
을 통해 (12)와 합성할 수 있다. 그리고 이 특수한 표현은 \(z=U_{i-1}(x)\)의 역변환에 따라 (28)에서 단순히 생긴다고 하자.

가정 (28)과 \(\xi,\zeta\)의 선형독립성에서, 위에서 (23)에 이어 쓴 것과 똑같은 논증으로 같은 대표원계를 대응시켜 \(\Gmod_{i-1}\mid\Mmod_{i-1}\)과 \(\Gmod^*_{i-1}\mid\Mmod^*_{i-1}\)의 동형을 얻는다. 마찬가지로 \(\Gmod^*_{i-1}\)이 \(\Mmod^*_{i-1}\)의 기초모듈이고, \(\Mmod_{i-1}\)이 \(1\)과 다른 초등인자를 유한개만 가지며 그것들이 \(\Mmod^*_{i-1}\)의 \(1\)과 다른 초등인자들임을 얻는다.

귀납단계를 수행하려면 먼저 \(\Gmod^*_{i-1}\mid\Mmod^*_{i-1}\), 따라서 \(\Gmod_{i-1}\mid\Mmod_{i-1}\)에 대하여 \(x_i\)의 차수가 제한된 대표원계를 택할 수 있음을 다시 보여야 한다. 이는 정리 III의 (7)에서 증명한 몫표현
\begin{equation}
\Gmod^*_{i-1}=\Mmod^*_{i-1}/(C^{(i)})
\tag{30}
\end{equation}
에서 따라온다. 여기서 \(C^{(i)}\)는 \(\Mmod^*_{i-1}\)에서 얻는 항등적으로 영이 아닌 임의의 \(\rho\)차 행렬식이고, \(\rho\)는 \(\Mmod^*_{i-1}\)의 랭크이다. 특수한 변환 (29)에 관한 가정에서 \(C^{(i)}\)를 언제나 \(x_i\)에 관하여 정칙으로 택할 수 있다는 것이 따라온다. 실제로 \(\Mmod^*_{i-1}\)에 대응하는 모듈 \(\widetilde\Mmod^*_{i-1}\)은 같은 랭크를 갖는다. \(\widetilde C^{(i)}\)가 \(\widetilde\Mmod^*_{i-1}\)에서 얻는 항등적으로 영이 아닌 임의의 \(\rho\)차 행렬식이고 \(z=U_{i-1}(x)\)에 따라 정칙인 \(C^{(i)}\)로 옮겨진다면, 가정에 따라 \(C^{(i)}\)는 \(\Mmod^*_{i-1}\)에서 얻는 \(\rho\)차 행렬식이다.

\(C^{(i)}\)가 존재하면 (8)에 대응하여 \(\xi\)들의 번호를 알맞게 정했을 때 특별한 모양의 \(\rho\)개 선형형식
\[
L_\lambda(\xi)
=C^{(i)}\xi_\lambda+c_{\lambda,1}^{(i)}\xi_{\rho+1}
+\cdots+c_{\lambda,s-\rho}^{(i)}\xi_s,
\qquad(\lambda=1,\ldots,\rho)
\]
이 \(\Mmod^*_{i-1}\)에 속한다. 이제 \(\xi\)의 모든 선형형식은
\begin{equation}
\begin{aligned}
H(\xi)={}&a_1^{(i)}\xi_1+\cdots+a_\rho^{(i)}\xi_\rho\\
&+b_1^{(i)}\xi_{\rho+1}+\cdots+b_{s-\rho}^{(i)}\xi_s
\pmod{L_1(\xi),\ldots,L_\rho(\xi)}
\end{aligned}
\tag{31}
\end{equation}
로 만들 수 있다. 여기서 \(a_1^{(i)},\ldots,a_\rho^{(i)}\)는 \(x_i\)의 차수가 제한되어 \(C^{(i)}\)의 차수 \(k\)에 미치지 않는다. 이 표현은 유일하다. 실제로 왼쪽이 \(L_1(\xi),\ldots,L_\rho(\xi)\)를 법으로 영이면, \(x_i\)의 차수를 고려하면서 \(\xi_1,\ldots,\xi_\rho\)의 계수를 비교하여 모든 \(a^{(i)},b^{(i)}\)가 항등적으로 사라짐을 얻는다.

특히
\[
H(\xi)\equiv0\pmod{\Gmod^*_{i-1}},
\qquad\text{따라서 (30)에 따라}\qquad
C^{(i)}H(\xi)\equiv0\pmod{\Mmod^*_{i-1}}
\]
라고 하자. 정리 III의 증명에서처럼
\[
C^{(i)}H(\xi)
-\sum_{\tau=1}^{s-\rho}
\left\{C^{(i)}b_\tau^{(i)}
-\sum_\lambda a_\lambda^{(i)}c_{\lambda,\tau}^{(i)}\right\}
\xi_{\rho+\tau}
\equiv0\pmod{\Mmod^*_{i-1}}
\]
을 얻는다. 정리 III에서와 같이 \(\xi_{\rho+1},\ldots,\xi_s\)의 계수들이 사라져야 하므로
\[
C^{(i)}b_\tau^{(i)}
=\sum_\lambda a_\lambda^{(i)}c_{\lambda,\tau}^{(i)},
\qquad(\tau=1,\ldots,s-\rho).
\]
따라서 \(b_\tau^{(i)}\)들의 차수도 제한되며 \(c_{\lambda,\tau}^{(i)}\)들의 최대차수에 미치지 않는다. 이로써 \(\Gmod^*_{i-1}\mid\Mmod^*_{i-1}\), 즉 \(\Gmod_{i-1}\mid\Mmod_{i-1}\)에 대하여 \(x_i\)의 차수가 어떤 고정된 차수 \(r\)에 미치지 않는 대표원계가 존재함이 증명되었다.

이제 유한개의 거듭제곱곱
\[
x_i^\alpha\xi_\lambda
\quad(\alpha=0,1,\ldots,k-1;\ \lambda=1,2,\ldots,\rho),
\]
\[
x_i^\beta\xi_{\rho+\tau}
\quad(\beta=0,1,\ldots,r-1;\ \tau=1,2,\ldots,s-\rho)
\]
를 \(\vartheta_1,\ldots,\vartheta_t\)라 하자. 또한 가산무한개의 식
\[
x_i^\gamma L_\lambda
\quad(\gamma=0,1,\ldots;\ \lambda=1,2,\ldots,\rho),
\qquad
x_i^\gamma\zeta_\nu
\quad(\gamma,\nu=0,1,\ldots)
\]
을 단순무한열 \(\omega_0,\omega_1,\ldots,\omega_\mu,\ldots\)로 배열하자. 그러면 \(\vartheta\)들과 \(\omega\)들은 정수적 원소 \(b^{(i+1)},c^{(i+1)},\ldots\)에 관하여 각 계 안에서도, 두 계 사이에서도 선형독립이다. 실제로 각 합이 유한개의 항에만 걸치는 관계
\[
\begin{aligned}
&\sum c_{\alpha\lambda}^{(i+1)}x_i^\alpha\xi_\lambda
+\sum c_{\beta\tau}^{(i+1)}x_i^\beta\xi_{\rho+\tau}\\
&\quad+\sum d_{\gamma\lambda}^{(i+1)}x_i^\gamma L_\lambda(\xi)
+\sum e_{\gamma\nu}^{(i+1)}x_i^\gamma\zeta_\nu=0
\end{aligned}
\]
에서, 가정한 \(\xi,\zeta\)의 선형독립성과 정수적 원소 \(b^{(i)}\)에 관한 \(\xi\)들의 상호 선형독립성 때문에 모든 \(e_{\gamma\nu}^{(i+1)}\)가 사라진다. 표현 (31)의 유일성에서 모든 \(c_{\alpha\lambda}^{(i+1)},c_{\beta\tau}^{(i+1)}\)가 사라지고, 끝으로 \(L_\lambda(\xi)\)들의 선형독립성에서 모든 \(d_{\gamma\lambda}^{(i+1)}\)가 사라진다.

이제 모듈 \((\Gmod_{i-1},L_1(\xi),\ldots,L_\rho(\xi))\)을 정수적 원소 \(b^{(i+1)}\)에 관한 모듈 \(\mathfrak C_i\)로 보자. \(\Gmod_{i-1},L_1(\xi),\ldots,L_\rho(\xi)\)가 \(\Mmod_{i-1}\)로 나누어지므로 \(\mathfrak C_i\)는 \(\Mmod_i\)로 나누어지고,
\[
\mathfrak C_i=(\omega_0,\omega_1,\ldots,\omega_\mu,\ldots)
\]
이다. 각 \(H(x)\equiv0\pmod{\gideal_{i-1}}\)에 대하여 (28), (31)과 방금 증명한 결과에서
\[
H(x)\equiv
b_1^{(i+1)}\vartheta_1+\cdots+b_t^{(i+1)}\vartheta_t
\pmod{\mathfrak C_i}
\]
라는, 정수적 원소 \(b^{(i+1)},c^{(i+1)}\)에 관한 모듈등식을 얻는다. 한편 \(\gideal_i\)는 \(\gideal_{i-1}\)로 나누어지고 \(\mideal\)은 \(\gideal_i\)로 나누어진다. 특히 \(\mathfrak C_i\)의 각 \(G(x)\)와 \(\Mmod_i\)의 각 \(F(x)\)에 대하여
\[
G(x)\equiv g_1^{(i+1)}\vartheta_1+\cdots+
g_t^{(i+1)}\vartheta_t\pmod{\mathfrak C_i},
\]
\[
F(x)\equiv a_1^{(i+1)}\vartheta_1+\cdots+
a_t^{(i+1)}\vartheta_t\pmod{\mathfrak C_i}
\]
로 두자. 모든 \(g(\vartheta)\)의 모듈을 \(\Gmod_i^*\), 모든 \(a(\vartheta)\)의 모듈을 \(\Mmod_i^*\)라 하면, (23)으로 이끈 것과 정확히 같은 논증으로
\begin{equation}
\Gmod_i=(\Gmod_i^*,\mathfrak C_i),\qquad
\Mmod_i=(\Mmod_i^*,\mathfrak C_i),\qquad
\mathfrak C_i=(\omega_0,\omega_1,\ldots,\omega_\mu,\ldots)
\tag{32}
\end{equation}
을 얻는다.

(32)를 유도할 때에도 \(C^{(i)}\)가 \(x_i\)에 관하여 정칙이라는 사실만을 썼다. 따라서 지표를 하나 올려 특수한 변환 (29)을 기초로 해도 논증 전체가 유지된다. (27)에 이어진 것과 정확히 같은 논증은 이 특수한 표현이 \(z=U_i(x)\)의 역변환에 따라 (32)에서 단순히 생김을 보여 준다.

이로써 일반적인 귀납단계의 가정 (28) 등이 한 단계 높은 지표에 대하여 모두 증명되었다. 앞에서 보았듯 이 가정들에서 정리 VII이 직접 따라오므로, 정리 VII은 일반적으로 증명되었다. 동시에 \(C^{(i)}\)들의 선택에 따른 \((\Gmod^*_{i-1},\Mmod^*_{i-1})\)의 임의성은 잉여류계 \(\Gmod^*_{i-1}\mid\Mmod^*_{i-1}\)의 동형을 통해 다시 사라진다는 것도 보였다.

특히 제3절에서 계수체로 삼은 \(P\)의 표수가 영, 즉 \(P\)에 포함되고 단위원으로 생성되는 소체가 유리수체의 유형이면, 각 \(i=1,\ldots,n\)에 대하여 \(C^{(i)}\)가 \(x_i\)에 관하여 정칙으로 남도록 부정원 \(u_{\mu\nu}\)들을 \(P\)의 원소 \(\bar u_{\mu\nu}\)로 특수화할 수 있다. 위의 논증은 이 특수화 아래에서도 정리 VII이 유지됨을 보여 준다\srcfn{14)}{헨첼은 임의의 \(P\) 대신 복소수체만을 택하고 주로 이 마지막 경우를 고찰했으며, 본래의 소거론에서만 부정원들을 기초로 삼았다. 그는 이어서, 본문에서 제시한 추론들을 본질적으로 사용하여 종결식형을 구성할 때 각 \(\Mmod_{i-1}\)에서 \(x\)에 관하여 어떤 고정된 경계를 넘는 유한한 차수까지만 진행하면 충분함을 보였다. 그러나 정리 VII이 주는 \(\Gmod_{i-1}\mid\Mmod_{i-1}\)과 \(\Gmod^*_{i-1}\mid\Mmod^*_{i-1}\)의 동형이 빠져 있어서, 차수들에 대한 이 독립성의 근거는 드러나지 않는다. (E. N.)}. 다만 기초모듈과 초등인자들은 일반적인 경우를 \(u_{\mu\nu}=\bar u_{\mu\nu}\)로 특수화하여 반드시 생길 필요는 없다\srcfn{15)}{\(\mideal=(x^2,ux+y)\)이면 \(\gideal_0=(1)\), \(\gideal_1=(1)\)이다. \(C^{(1)}=x^2\)을 택하면 \(\Gmod_1^*=(1,x)\), \(\Mmod_1^*=(ux+y,xy)\)이고, \(\Mmod_1^*\)의 초등인자는 \(y^2,1\)이며 \(C^{(2)}=y^2\)으로 택할 수 있다. \(u=0\)일 때에도 \(C^{(1)},C^{(2)}\)은 각각 \(x,y\)에 관하여 정칙으로 남지만, \(\Mmod_1^*=(y,xy)\)의 초등인자는 \(y,y\)이다. 따라서 이 경우에도 \(\Gmod_1\mid\Mmod_1\)은 유한모듈의 잉여류계와 동형이지만 일반적인 경우의 잉여류계와 동형은 아니다. 반대로 \(C^{(1)}=ux+y\)를 택하고 \(C^{(1)}\)이 정칙으로 남도록 특수화했다면 동형도 유지되었을 것이다. (E. N.)}.

\section*{\S{}5. 다항식 아이디얼의 종결식형과 초등인자형}

\(x_{i+1},\ldots,x_n\)을 \(P(u)\)에 첨가하자. 그러면 \(\Gmod_{i-1},\Mmod_{i-1}\), 따라서 \(\Gmod^*_{i-1},\Mmod^*_{i-1}\)도 정수적 원소들을 한 부정원 \(x_i\)의 다항식으로 볼 수 있는 선형형식 모듈이 된다. 정리 VII에 따라 이 경우 \(\Mmod_{i-1}\)의 \(1\)과 다른 초등인자들과 기껏해야 유한개의 초등인자 \(1\)은 \(\Mmod^*_{i-1}\)의 초등인자들과 일치한다. 이 초등인자들의 곱, 즉 \(\Mmod^*_{i-1}\)의 \(\rho\)번째 행렬식인자는 \(\Gmod^*_{i-1}\)에서 \(\Mmod^*_{i-1}\)로 넘어가는 치환의 행렬식과 같으며, 정의 III에 따라 \(\Norm(\Gmod^*_{i-1}\mid\Mmod^*_{i-1})\)이다. (24), 또는 일반적인 \(i\)에 대하여 (28)에서 생기는 유사한 공식은 이 초등인자들의 곱이 \(\Gmod_{i-1}\)에서 \(\Mmod_{i-1}\)로 넘어가는 치환의 행렬식과도 같음을 보여 준다. 따라서 이를 \(\Norm(\Gmod_{i-1}\mid\Mmod_{i-1})\)이라 써야 한다. 즉
\[
\Norm(\Gmod_{i-1}\mid\Mmod_{i-1})
=\Norm(\Gmod^*_{i-1}\mid\Mmod^*_{i-1})
=R^{(i)}(x).
\]
다항식 \(R^{(i)}(x)\)은 \(\Mmod^*_{i-1}\)에서 얻는 모든 \(\rho\)차 행렬식의 다항식 의미의 최대공약원이다. 이를 \(x_{i+1},\ldots,x_n\)에 관하여 정수적이고 원시적인 다항식이라고 가정할 수 있으므로, \(C^{(i)}\)의 인자로서 \(x_i\)에 관하여 정칙이다. 마찬가지로 \(\Mmod_{i-1}\), 또는 \(\Mmod^*_{i-1}\)의 가장 높은 초등인자 \(E^{(i)}(x)\)도 \(x_{i+1},\ldots,x_n\)에 관하여 정수적이고 원시적이며, 따라서 \(x_i\)에 관하여 정칙이라고 가정할 수 있다. 이로써 다음 정의에 이른다.

\medskip
\noindent\textbf{정의 VI.} \emph{다항식 \(R^{(i)}(x)\)을 \(\mideal\)의 \(i\)단계 종결식이라 하고 \(E^{(i)}(x)\)을 \(i\)단계 초등인자라 한다. 따라서 \(i\)단계 종결식은 기초아이디얼 \(\gideal_{i-1}\)의 \(\mideal\)에 대한 노름, 즉 \(x_{i+1},\ldots,x_n\)을 \(P(u)\)에 첨가했을 때 모듈 \(\Gmod_{i-1}\)의 \(\Mmod_{i-1}\)에 대한 노름이다. 같은 첨가 아래 \(E^{(i)}(x)\)은 몫 \(\Mmod_{i-1}/\Gmod_{i-1}\)과 같다. 곱
\[
R_\mideal=R^{(1)}R^{(2)}\cdots R^{(n)},\qquad
E_\mideal=E^{(1)}E^{(2)}\cdots E^{(n)}
\]
을 각각 \(\mideal\)의 종결식형과 초등인자형이라 한다.}

종결식형과 초등인자형에 대해서는 다음이 성립한다.

\medskip
\noindent\textbf{정리 VIII.} \emph{종결식형은 초등인자형으로 나누어지고, 거꾸로 \(E_\mideal\)의 어떤 거듭제곱은 \(R_\mideal\)로 나누어진다. \(E_\mideal,R_\mideal\)은 \(\mideal\)로 나누어진다. 더 일반적으로
\[
E^{(i)}\cdots E^{(n)}\gideal_{i-1}\zero{\mideal},
\]
따라서
\[
R^{(i)}\cdots R^{(n)}\gideal_{i-1}\zero{\mideal}.
\]}

노름은 가장 높은 초등인자로 나누어지고, 거꾸로 이 초등인자의 어떤 거듭제곱은 노름으로 나누어진다. 따라서 \(x_{i+1},\ldots,x_n\)을 첨가했을 때 \(R^{(i)}(x),E^{(i)}(x)\)에 대하여 이 나눗셈이 성립한다. 그런데 \(R^{(i)}(x),E^{(i)}(x)\)을 \(x_{i+1},\ldots,x_n\)에 관한 원시다항식으로 가정했으므로, 모든 부정원 \(x_i,x_{i+1},\ldots,x_n\)에 관하여도 이 나눗셈이 성립한다. 따라서 이 표현들의 정의에 따라 모든 부정원 \(x\)에 관하여 \(R_\mideal\)은 \(E_\mideal\)로 나누어지고 \(E_\mideal\)의 어떤 거듭제곱은 \(R_\mideal\)로 나누어진다.

정리 VII과 정리 II에 따라 \(x_{i+1},\ldots,x_n\)을 첨가했을 때 가장 높은 초등인자 \(E^{(i)}(x)\)은 몫 \(\Mmod_{i-1}/\Gmod_{i-1}\)으로 정의된다. 따라서 다시 모든 부정원 \(x\)의 다항식으로 돌아오면 각
\begin{equation}
G(x)\zero{\gideal_{i-1}}
\quad\text{에 대하여 어떤}\quad b^{(i+1)}\ne0
\quad\text{가 존재하여}\quad
b^{(i+1)}E^{(i)}(x)G(x)\zero{\mideal}
\tag{33}
\end{equation}
이 된다. 특히 \(\gideal_0\)은 단위원아이디얼과 같으므로
\[
b^{(2)}E^{(1)}(x)\zero{\mideal},\quad b^{(2)}\ne0,
\qquad\text{따라서}\qquad
E^{(1)}(x)\zero{\gideal_1}.
\]
일반적으로 (33)에 따라
\[
E^{(1)}(x)\cdots E^{(i-1)}(x)\zero{\gideal_{i-1}}
\]
이면 어떤 \(b^{(i+1)}\ne0\)가 존재하여
\[
b^{(i+1)}E^{(1)}\cdots E^{(i)}\zero{\mideal}
\]
이고, 따라서
\[
E^{(1)}\cdots E^{(i)}\zero{\gideal_i}.
\]
\(\gideal_n=\mideal\)이므로
\begin{equation}
E_\mideal=E^{(1)}\cdots E^{(n)}\zero{\mideal},
\qquad\text{따라서}\qquad
R_\mideal=R^{(1)}\cdots R^{(n)}\zero{\mideal}.
\tag{34}
\end{equation}

(33)의 \(G(x)\)가 \(\gideal_{i-1}\)의 한 아이디얼기저에 속하는 유한개의 다항식을 돌게 하고, \(c^{(i+1)}(x)\)를 이에 대응하는 \(b^{(i+1)}(x)\)들의 곱이라고 하자. 그러면
\[
c^{(i+1)}E^{(i)}(x)\gideal_{i-1}\zero{\mideal},
\qquad c^{(i+1)}\ne0,
\]
따라서
\[
E^{(i)}(x)\gideal_{i-1}\zero{\gideal_i}.
\]
이를 위와 같이 유한번 반복하면
\[
E^{(n)}(x)\cdots E^{(i)}(x)\gideal_{i-1}\zero{\mideal}
\]
이고, 따라서
\[
R^{(n)}(x)\cdots R^{(i)}(x)\gideal_{i-1}\zero{\mideal}.
\]
이로써 정리 VIII의 모든 부분이 증명되었다.

정리 VIII은 본질적으로 초등인자형의 성질에 바탕을 둔다. 그러나 종결식형이 \(\mideal\)에 대하여 특징적인 의미를 갖는다는 것은 다음 정리가 보여 준다.

\medskip
\noindent\textbf{정리 IX.} \emph{\(\nideal\)이 아이디얼의 의미에서 \(\mideal\)의 인자이고 종결식형 \(R_\nideal,R_\mideal\)이 같으면 아이디얼 \(\nideal,\mideal\)도 같다.}

\(\gideal_0,\ldots,\gideal_{n-1},\gideal_n=\mideal\)과 \(\hideal_0,\ldots,\hideal_{n-1},\hideal_n=\nideal\)이 각각 \(\mideal,\nideal\)의 기초아이디얼들이라고 하자. 먼저 \(\gideal_0,\hideal_0\)은 둘 다 단위원아이디얼이므로 \(\gideal_0=\hideal_0\)이다. 이제 \(\gideal_{i-1}=\hideal_{i-1}\)이라고 가정하자. 그러면 \(\Mmod_{i-1},\Nmod_{i-1}\)은 같은 기초모듈 \(\Gmod_{i-1}\)을 갖는다. 가정 \(R_\nideal=R_\mideal\)은 \(x_{i+1},\ldots,x_n\)을 첨가했을 때
\[
\Norm(\Gmod_{i-1}\mid\Nmod_{i-1})
=\Norm(\Gmod_{i-1}\mid\Mmod_{i-1}),
\]
또는
\[
\Norm(\Gmod^*_{i-1}\mid\Nmod^*_{i-1})
=\Norm(\Gmod^*_{i-1}\mid\Mmod^*_{i-1})
\]
이라고 쓸 수 있다. 여기서 (32)에 대응하여
\[
\Nmod_{i-1}=(\Nmod^*_{i-1},\mathfrak C_{i-1})
\]
로 두었다. 따라서 \(\Nmod^*_{i-1}\)도 \(\xi\)의 선형형식 모듈이고 \(\Gmod^*_{i-1}\)은 그 기초모듈이다. \(\Mmod_{i-1}\)이 \(\Nmod_{i-1}\)로 나누어지므로 \(\Mmod^*_{i-1}\)도 \(\Nmod^*_{i-1}\)로 나누어진다. 따라서 정리 IV에 따라 이 첨가 아래 두 모듈 \(\Mmod^*_{i-1},\Nmod^*_{i-1}\), 그러므로 \(\Nmod_{i-1},\Mmod_{i-1}\)도 일치한다. 한편 \(x_{i+1},\ldots,x_n\)의 첨가는 \(\Mmod_{i-1}\), 또는 \(\mideal\)에
\[
b^{(i+1)}(x)G(x)\zero{\mideal},\qquad b^{(i+1)}\ne0
\]
이 되는 모든 다항식 \(G(x)\)를 더하는 것을 뜻한다. 즉 이 첨가로 \(\Mmod_{i-1}\), 또는 \(\mideal\)은 \(\gideal_i\)로 옮겨지고, 마찬가지로 \(\nideal\)은 \(\hideal_i\)로 옮겨진다. 따라서 \(\gideal_i=\hideal_i\)이고, 끝내 \(\gideal_n=\hideal_n\), 즉 \(\mideal=\nideal\)이다.

끝으로 같은 이행원리를 정리 V에 적용하면 다음을 얻는다.

\medskip
\noindent\textbf{정리 X.} \emph{\(R^{(i)}=S^{(i)}T^{(i)}\)가 \(R^{(i)}\)을 다항식의 의미에서 서로소인 두 다항식 \(S^{(i)},T^{(i)}\)으로 분해한 것이라고 하자. 그러면 \(\mideal\)의 인자이고 같은 \((i-1)\)단계 기초아이디얼 \(\gideal_{i-1}\)을 갖는 아이디얼 \(\mathfrak s,\tideal\)이 각각 유일하게 존재하여
\[
S^{(i)}=\Norm(\Gmod_{i-1}\mid\Smod_{i-1}),\qquad
T^{(i)}=\Norm(\Gmod_{i-1}\mid\Tmod_{i-1})
\]
이 된다. \(\mathfrak s,\tideal\)은 각각
\[
s^{(i+1)}(x)T^{(i)}(x)S(x)\zero{\mideal},
\qquad s^{(i+1)}\ne0,
\]
\[
t^{(i+1)}(x)S^{(i)}(x)T(x)\zero{\mideal},
\qquad t^{(i+1)}\ne0
\]
이 되는 다항식 \(S(x)\), 또는 \(T(x)\) 전체로 정의된다. 또한 \(\gideal_i\)는 \(\mathfrak s,\tideal\)의 최소공배아이디얼이다.
\[
\gideal_i=[\mathfrak s,\tideal].
\]}

여기서는 \(\mathfrak s,\tideal\)의 기저원소들에 대응하는 \(s^{(i+1)},t^{(i+1)}\)들의 곱을 취하여, \(s^{(i+1)},t^{(i+1)}\)을 각 아이디얼에 대하여 고정해도 된다는 점만 덧붙이면 된다. 또한 앞에서 보았듯 \(x_{i+1},\ldots,x_n\)을 첨가하면 \(\mideal\)은 \(\gideal_i\)로 바뀐다. 특히 \(R^{(i)}\)의 분해를 기약다항식들의 거듭제곱, 즉 으뜸인자들까지 계속할 수 있으며, 그 결과 \(\gideal_i\)도 이에 대응하는 최소공배아이디얼 표현을 갖는다.

\section*{\S{}6. 본래의 소거론}

정리 VIII, 즉 식 (34)에 따르면 종결식형과 초등인자형은 \(\mideal\)로 나누어지므로 \(\mideal\)의 모든 영점에서 사라진다. 여기서 \glqq{}\(\mideal\)의 영점\grqq{}이란 다음 값의 계를 뜻한다. \(x_{i+1},\ldots,x_n\)을 계수영역에 첨가했다고 생각하고, \(P(u;x_{i+1},\ldots,x_n)\)에서 유도되는 대수적으로 닫힌 체에 속하는 \(x_1,\ldots,x_i\)의 값의 계 가운데 \(\mideal\)의 모든 다항식이 사라지는 것들이다. 이는 \(i=1,\ldots,n\) 각각에 대하여 말한다\srcfn{16)}{모든 체가 본질적으로 유일한 방식으로 대수적으로 닫힌 체까지 확대될 수 있다는 사실은 Steinitz, \emph{Algebraische Theorie der Körper}, \emph{J. f. M.} 137 (1910), 167--309쪽에서 증명되었고, 이로써 대수학의 기본정리에 대한 유리적 등가물이 주어졌다. 실제로 각 \(i=1,\ldots,n\)에 대해서는 \(P(u,x_{i+1},\ldots,x_n)\)의 유한확대체만 다룬다. (E. N.)}. 뒤에서 보이듯 첨가한 \(x_{i+1},\ldots,x_n\)을 \(P(u)\)의 원소들로 바꿀 수도 있다. 거꾸로 종결식형의 영점에서 \(\mideal\)의 영점들을 유도하는 것이 소거론의 내용이다. 이 역방향은 이 절에서 전개할 연속소거에 바탕을 둔다. 다음 절에서 그 과정에 나타나는 형식 \(D^{(i)}(x)\)이 \(E^{(i)}(x)\)으로 나누어짐을 보이면, \(E^{(i)}(x)\)의 어떤 거듭제곱이 \(R^{(i)}(x)\)으로 나누어진다는 사실에서 구하는 역방향이 따라온다.

여기서 제시할 연속소거 이론은 다음 정리에 바탕을 둔다.

\medskip
\noindent\textbf{정리 XI.} \emph{\(\bideal\)이 \(P\)를 계수체로 갖는 한 부정원 \(t\)의 다항식 아이디얼, 따라서 주아이디얼이라고 하자. \(h(t)\)는 \(\bideal\)에 속하는 \(k\)차 다항식이고, \(\Bmod\)은 \(\bideal\)을 \(h(t)\)로 법을 취했을 때 생기는 \(1,t,\ldots,t^{k-1}\)의 선형형식 모듈이라고 하자. \(\bideal\)의 기저다항식 \(f(t)\)의 차수가 \(p\)이면 \(\Bmod\)의 랭크는 \(k-p\)이다.}

가정에 따라 \(h(t)=h_1(t)f(t)\)이고 \(h_1(t)\)의 차수는 \(k-p\)이다. 따라서 다항식
\[
f(t),tf(t),\ldots,t^{k-p-1}f(t)
\]
이 대표하는 \(\bideal\)의 \(h(t)\)에 관한 잉여류들은 선형독립이다. 또한 \(\bideal\)의 각 잉여류는 \(P\)의 계수를 갖는 선형결합으로 이 \(k-p\)개의 특별한 류들로 표현된다. 한편 \(\bideal\)은 \(h(t)\)를 법으로 하여 \(1,t,\ldots,t^{k-1}\)의 선형형식 모듈로 옮겨지므로 그 랭크는 정확히 \(k-p\)이다.

이제 \(\aideal=\aideal_1\)이 변환된 다항식 아이디얼이고, \(A^{(1)}(x)\)가 \(\aideal_1\)에 속하면서 \(x_1\)에 관하여 정칙인 \(k_1\)차 다항식이라고 하자. \(\mathfrak A_1\)은 \(\aideal_1\)을 \(A^{(1)}(x)\)로 법을 취했을 때 생기는 \(1,x_1,\ldots,x_1^{k_1-1}\)의 선형형식 모듈이고, 그 계수는 다항식 \(a^{(2)}(x)\)이다. 또 \(\aideal_2\)는 \(\mathfrak A_1\)에서 얻는 모든 \(k_1\)차 행렬식이 생성하는 \(x_2,\ldots,x_n\)의 아이디얼이라고 하자. 정리 XI에 따라 \(\aideal_2\)가 영아이디얼인 것은 \(\aideal_1\)의 다항식들이 \(x_1\)에 관하여 차수 \(p_1>0\)인 공통인자, 여기서는 다항식 의미의 인자를 가질 때 그리고 그때뿐이다. \(\aideal_2\)가 영아이디얼이 아니면, \(\aideal\)이 변환된 아이디얼이라는 가정 때문에 \(\aideal_2\)는 \(x_2\)에 관하여 정칙인 \(k_2\)차 다항식 \(A^{(2)}(x)\)을 포함한다. 이는 변환 (12)을 특수한 변환 (25), (26)에서 합성하면 직접 알 수 있다. 다시 \(\mathfrak A_2\)를 \(\aideal_2\)가 \(A^{(2)}(x)\)을 법으로 하여 옮겨지는 \(1,x_2,\ldots,x_2^{k_2-1}\)의 선형형식 모듈이라 하고, \(\aideal_3\)을 \(\mathfrak A_2\)의 모든 \(k_2\)차 행렬식이 생성하는 \(x_3,\ldots,x_n\)의 아이디얼이라 하자. 이 아이디얼은 \(x_3\)에 관하여 정칙인 다항식 \(A^{(3)}(x)\)을 포함해야 한다.

이 절차를 영아이디얼에 이르거나 \(\aideal_{n+1}\)이 단위원아이디얼이 될 때까지 계속한다. \(\aideal_{n+1}\)은 \(x\)와 무관하므로 영아이디얼 또는 단위원아이디얼일 수밖에 없다. 아이디얼 \(\aideal_1,\aideal_2,\ldots\)은 기초로 택한 다항식 \(A^{(\lambda)}(x)\)과 무관하게 \(\aideal\)만으로 유일하게 정해진다. \(\aideal_1=\aideal\)에 대해서는 분명하므로 \(\aideal_1,\ldots,\aideal_i\)까지 성립한다고 가정하자. \(\aideal_i\)가 \(A^{(i)}(x)\)을 법으로 하여 선형형식 모듈 \(\mathfrak A_i\)로 옮겨진다면, \(\mathfrak A_i\)는 \(x_i\)의 차수가 \(k_i\)에 미치지 않는 \(\aideal_i\)의 모든 다항식으로도 특징지을 수 있다. 따라서 \(\mathfrak A_i\)는 \(A^{(i)}(x)\)의 차수 \(k_i\)에만 의존한다. 이제 \(\bar A^{(i)}(x)\)가 \(\aideal_i\)에 속하면서 \(x_i\)에 관하여 정칙인 \(\bar k_i>k_i\)차 다항식이고, \(\overline{\mathfrak A}_i\)가 대응하는 선형형식 모듈, \(\bar\aideal_{i+1}\)이 \(\overline{\mathfrak A}_i\)의 \(\bar k_i\)차 행렬식들이 생성하는 아이디얼이라고 하자. 그러면 모듈등식
\[
\overline{\mathfrak A}_i=
(\mathfrak A_i,A^{(i)},x_iA^{(i)},\ldots,
x_i^{\bar k_i-k_i-1}A^{(i)})
\]
이 성립한다. \(A^{(i)}(x)\)이 \(x_i\)에 관하여 정칙이므로 \(A^{(i)},x_iA^{(i)},\ldots\)를 선형형식의 새 부정원들로 도입할 수 있다. 따라서 \(\mathfrak A_i\)의 \(k_i\)차 행렬식들과 \(\overline{\mathfrak A}_i\)의 \(\bar k_i\)차 행렬식들이 일치하고, \(\bar\aideal_{i+1}=\aideal_{i+1}\)이다\srcfn{17)}{이는 원리상 \(\Gmod^*_{i-1}\mid\Mmod^*_{i-1}\)과 \(\Gmod_{i-1}\mid\Mmod_{i-1}\)의 동형이 바탕을 둔 것과 같은 추론이다.}.

또한 \(\aideal_{i+1}\)은 언제나 \(\aideal_i\)로 나누어진다. \(\aideal_{i+1}\)이 영아이디얼이면 분명하다. 그렇지 않으면 \(\mathfrak A_i\)의 랭크는 \(k_i\)이다. 따라서 정의상 \(\aideal_{i+1}\)은 \(\mathfrak A_i\)의 \(k_i\)차 행렬식들로 이루어진 아이디얼이고, \(\mathfrak A_i\), 따라서 \(\aideal_i\)로 나누어진다. 그러므로 모든 \(\aideal_{i+1}\)은 \(\aideal\)로 나누어진다. 따라서 \(\aideal_{n+1}\)이 단위원아이디얼이면 \(\aideal\)도 단위원아이디얼이다.

이제 \(\aideal\)은 단위원아이디얼과 다르고
\[
\aideal_1\ne0,\ldots,\aideal_i\ne0,\qquad
\aideal_{i+1}=0
\]
이라고 하자. \(D^{(i)}(x)\)는 정리 XI에 따라 존재하는 \(\aideal_i\)의 모든 다항식의 최대공약원, 즉 다항식 의미의 공통인자이다. \(A^{(i)}(x)\)의 인자이므로 \(D^{(i)}(x)\) 자체도 \(x_i\)에 관하여 정칙이고 차수는 \(p_i>0\)이다. \(x_{i+1},\ldots,x_n\)을 \(P(u)\)에 첨가하자. 그러면 \(D^{(i)}(x)\)은 \(P(u;x_{i+1},\ldots,x_n)\)의 한 대수적 확대체에서 \(p_i\)개의 일차인자로 분해된다. \(x_i-\bar x_i\)가 그 한 인자라 하자. 그러면 \(x_i=\bar x_i\)에서 \(\aideal_i\)의 모든 다항식이 사라지고, 정리 XI에 따라 \(\aideal_{i-1}\)의 다항식들은 최대공약원 \(D^{(i-1)}(x)\)을 얻는다. \(A^{(i-1)}(x)\)의 인자이므로 \(D^{(i-1)}(x)\)도 \(x_{i-1}\)에 관하여 정칙이고 차수는 \(p_{i-1}>0\)이다. 특히 항등적으로 사라질 수 없고, \(P(u,\bar x_i,x_{i+1},\ldots,x_n)\)의 한 확대체에서 \(p_{i-1}\)개의 일차인자로 분해된다. \(x_{i-1}-\bar x_{i-1}\)가 그 한 인자이면 이에 \(\aideal_{i-2}\)의 최대공약원이 대응한다. 이 과정을 계속하면 \(P(u;x_{i+1},\ldots,x_n)\)의 한 대수적 확대체에 속하는 \(\aideal\)의 공통영점들을 유한개 얻는다. 거꾸로 \(\aideal_i\)는 \(\aideal\)로 나누어지므로 \(\aideal\)의 모든 영점은 \(\aideal_i\)의 영점이기도 하다. 특히
\[
x_1=\bar x_1,\ldots,x_i=\bar x_i,\quad
x_{i+1}=x_{i+1},\ldots,x_n=x_n
\]
이 \(\aideal\)의 영점이면 \(\aideal_{i+1}=0\)이고, \(\aideal_i\)의 다항식들은 일차인자 \(x_i-\bar x_i\)를 갖는 최대공약원 \(D^{(i)}(x)\)을 얻는다. 요약하면 다음과 같다.

\medskip
\noindent\textbf{정리 XII.} \emph{\(\aideal\)이 단위원아이디얼과 다르고
\[
\aideal_1\ne0,\ldots,\aideal_i\ne0,\qquad
\aideal_{i+1}=0
\]
이라고 하자. \(x_{i+1},\ldots,x_n\)을 \(P(u)\)에 첨가하면 \(P(u,x_{i+1},\ldots,x_n)\)의 한 대수적 확대체에 속하는 서로 결부된 값의 계 \(\bar x_1,\ldots,\bar x_i\)가 유한개 존재하여, \(x_1=\bar x_1,\ldots,x_i=\bar x_i\)에서 \(\aideal\)의 모든 다항식이 사라진다. 거꾸로 \(\aideal\)의 그러한 영점이 있으면 \(\aideal_{i+1}=0\)이고, \(\aideal_i\)의 모든 다항식은 일차인자 \(x_i-\bar x_i\)를 갖는 공통인자 \(D^{(i)}(x)\)를 갖는다.}

특히 정리 XII는 영점을 갖지 않는 모든 아이디얼 \(\aideal\)이 단위원아이디얼임을 보여 준다.

값의 계가 서로 결부되는 방식까지 명시적으로 보여 주는 분해를 얻기 위해, \(P(u,x_{i+1},\ldots,x_n)\)에 첨가할 새 부정원 \(v\)들을 도입한다. 다음과 같이 두자.
\begin{equation}
x_i=-v_1x_1-\cdots-v_{i-1}x_{i-1}+z_i.
\tag{35}
\end{equation}
(12)와 (35)를 합성하면 다시 (12)형의 치환을 얻는다. 다만 \(u_{i\kappa}\)는 \(w_{i\kappa}=u_{i\kappa}-v_\kappa\)로, 더 일반적으로 \(u_{i+\lambda,\kappa}\)는
\[
w_{i+\lambda,\kappa}
=u_{i+\lambda,\kappa}-u_{i+\lambda,i}v_\kappa
\]
로 바뀐다. 따라서 치환 (35)가 변환된 아이디얼 \(\aideal\)을 \(\bideal\)로 옮기면, \(\bideal\)은 \(\aideal\)에서 \(u_{\mu\nu}\)를 \(w_{\mu\nu}\)로 바꾸고 \(v\)들을 첨가하여 생긴다. 같은 과정으로 \(\bideal\)이 정의하는 아이디얼 \(\bideal_1,\bideal_2,\ldots\)은 \(\aideal_1,\aideal_2,\ldots\)에서 생긴다. 거꾸로 \(v=0\)으로 두면 \(\bideal_i\)에서 \(\aideal_i\)를 얻는다. 따라서 \(\aideal_i,\bideal_i\)는 동시에 영아이디얼이거나 동시에 영아이디얼과 다르다.

다시
\[
\aideal_1\ne0,\ldots,\aideal_i\ne0,\qquad\aideal_{i+1}=0
\]
이라 하자. 그러면 \(\bideal_1\ne0,\ldots,\bideal_i\ne0,\ \bideal_{i+1}=0\)도 성립한다. \(\bar x_1,\ldots,\bar x_i,x_{i+1},\ldots,x_n\)이 서로 결부된 \(\aideal\)의 영점계이고
\[
\bar z_i=\bar x_i+v_1\bar x_1+\cdots+v_{i-1}\bar x_{i-1}
\]
로 정의하면, (35)에 따라 \(\bar x_1,\ldots,\bar x_{i-1},\bar z_i,x_{i+1},\ldots,x_n\)은 \(\bideal\)의 영점이다. 이제 \(H^{(i)}(z,x)\)가 \(\bideal_i\)의 모든 다항식의 최대공약원이라고 하자. 이는 \(v\)에 관하여 정수적이고 원시적이라고 가정할 수 있다. 그러면 정리 XII의 마지막 부분에 따라 \(H^{(i)}(z,x)\)은 인자
\[
z_i-\bar z_i
=z_i-(\bar x_i+v_1\bar x_1+\cdots+v_{i-1}\bar x_{i-1})
\]
를 갖는다. \(x_{i+1},\ldots,x_n\)을 첨가했을 때 나타나는 \(\aideal\)의 각 영점에 그러한 인자가 대응한다.

이 인자들이 \(H^{(i)}(z,x)\)의 분해를 모두 소진함을 보이자. 즉 \(z\)와 \(v\)에 관하여 일차인자들로 분해되지 않는 인자 \(H_1^{(i)}\)를 \(H^{(i)}\)에서 떼어 낼 수 없음을 보이면 된다. 그렇다고 가정하자. \(v\)에 관한 원시성 때문에 \(H^{(i)}\)은 적어도 하나의 일차인자 \(z_i-\bar z_i\)를 가지며, \(\bar z_i\)는 \(P(u,v,x_{i+1},\ldots,x_n)\)의 한 대수적 확대체에 속한다. \(\bar x_1,\ldots,\bar x_{i-1},\bar z_i,x_{i+1},\ldots,x_n\)이 이에 대응하는 \(\bideal\)의 영점이면, 이는 \(x_{i+1},\ldots,x_n\)을 첨가했을 때 나타나는 유한개의 \(\aideal\)의 영점 가운데 하나와 일치해야 하므로 \(v\)와 무관하다. 따라서 \(\bar z_i\)는 \(v\)들을 반드시 선형으로 포함하며, 대수적으로나 유리비선형으로 포함하지 않는다. 그러면 가정에 어긋나게 \(H_1^{(i)}\)에서 \(z,v\)에 관한 또 하나의 일차인자
\[
z_i-(\bar x_i+v_1\bar x_1+\cdots+v_{i-1}\bar x_{i-1})
\]
를 떼어 낼 수 있다. 따라서 명시적인 분해
\[
H^{(i)}(z,x)=
\prod_\lambda
\left\{z_i-(\bar x_{i\lambda}+v_1\bar x_{1\lambda}
+\cdots+v_{i-1}\bar x_{i-1,\lambda})\right\}^{\alpha_\lambda}
\]
를 얻는다.

\(H^{(i)}\)의 차수와 각 지수 \(\alpha_\lambda\)는 \(D^{(i)}(x)\)의 대응하는 것들과 같아야 한다. \(H^{(i)}\)은 \(D^{(i)}\)에서 \(u_{\mu\nu}=w_{\mu\nu}\)로 특수화하여 생기고, \(D^{(i)}\)은 \(H^{(i)}\)에서 \(v=0\)으로 특수화하여 생기기 때문이다. 그러므로 이 분해는 또 다음을 보여 준다. 부정원 \(u_{\mu\nu}\)들을 \(P\)에 첨가했기 때문에, \(\aideal\)에서 시작하는 소거에서도 \(D^{(i)}\)의 각 영점 \(\bar x_i\)에는 \(\aideal_{i-1},\ldots,\aideal_1\)의 매 단계에서 일차인 최대공약원 또는 그 거듭제곱이 대응한다. 따라서 매번 \(\aideal\)의 서로 결부된 영점계 \(\bar x_1,\ldots,\bar x_i,x_{i+1},\ldots,x_n\) 하나만이 생긴다\srcfn{18)}{여기서 제시한 연속소거는 크로네커의 방법과 달리 주어진 아이디얼 \(\aideal\)에만 의존하고 어떤 아이디얼기저에도 의존하지 않는다. 따라서 특히 지수 \(\alpha_\lambda\)들도 기저와 무관하다. 헨첼에 따르면 이 방법으로 소거를 완전히 수행하려면 몫 \(\aideal/D^{(i)}\)에 대하여 절차를 계속해야 하지만, 다음 절을 고려하여 여기서는 생략한다. (E. N.)}.

\section*{\S{}7. 종결식형과 소거론}

종결식형과 소거론의 관계를 얻기 위해 제6절의 연속소거를 아이디얼 \(\mideal\)을 그 \((i-1)\)단계 기초아이디얼로 나눈 몫
\[
\aideal=\mideal/\gideal_{i-1}
\]
에 적용한다. 먼저 이 몫이 변환된 아이디얼임을 보여야 한다. \(\mideal\)은 변환되었고, 따라서 제3절의 정리 VI에 따라 \(\gideal_{i-1}\)도 변환되었다. 이제
\[
H(x)\zero{\mideal/\gideal_{i-1}},\qquad
H(x)=\overline H(y)=\sum_i\alpha_i(u)\overline H_i(y)
=\sum_i\alpha_i(u)H_i(x)
\]
이면
\[
\overline H(y)\,\overline\gideal_{i-1,(u)}
\zero{\overline\mideal_{(u)}},
\qquad\text{따라서}\qquad
\overline H(y)\,\overline\gideal_{i-1}
\zero{\overline\mideal_{(u)}}.
\]
그러므로
\[
\overline H_i(y)\,\overline\gideal_{i-1}
\zero{\overline\mideal},
\qquad\text{따라서}\qquad
H_i(x)\gideal_{i-1}\zero{\mideal}.
\]
즉 \(H_i(x)\)는 \(\mideal/\gideal_{i-1}\)로 나누어지고, 이 아이디얼이 변환되었음이 증명되었다. 따라서 제6절의 소거법을 적용할 수 있다.

한편 제4절의 (28)을 고려하면 (30)은 \(C^{(i)}(x)\)가 \(\mideal/\gideal_{i-1}\)로 나누어짐을 보여 준다. 여기서 \(C^{(i)}(x)\)는 \(\Mmod^*_{i-1}\)에서 얻는 항등적으로 영이 아닌 임의의 \(\rho\)차 행렬식이다. \(i>1\)이면 \(C^{(i)}\)은 \(x_1\)에 관하여 0차이므로, 아이디얼 \(\aideal=\aideal_1\)에 대응하는 모듈 \(\mathfrak A_1\)은
\[
C^{(i)},x_1C^{(i)},\ldots,x_1^{k_1-1}C^{(i)}
\]
을 포함한다. 따라서 \((C^{(i)})^{k_1}\)은 \(\aideal_2\)로 나누어진다. 마찬가지로 \((C^{(i)})^{k_1k_2}\)은 \(\aideal_3\)으로 나누어지고, 끝으로
\[
(C^{(i)})^{k_1k_2\cdots k_{i-1}}
\]
은 \(\aideal_i\)로 나누어진다. 그러므로 \(\aideal_1,\ldots,\aideal_i\)는 영아이디얼과 다르며, 이는 \(i=1\)일 때에도 성립한다. 이제 \(D^{(i)}(x)\)가 \(\aideal_i\)의 모든 다항식의 최대공약원이라고 하자. 그 차수가 \(p_i>0\)인 것은 \(\aideal_{i+1}\)이 영아이디얼일 때뿐이다. \(D^{(i)}(x)\)의 정의와 \(\aideal_i\)가 \(\aideal\)로 나누어진다는 사실에서
\[
d^{(i+1)}D^{(i)}(x)\zero{\mideal/\gideal_{i-1}},
\qquad d^{(i+1)}\ne0
\]
을 얻는다. 따라서 \(d^{(i+1)}D^{(i)}\)은 \(x_1,\ldots,x_{i-1}\)과 무관하고 \(\mideal/\gideal_{i-1}\)에 속하는 다항식이다. 한편 \(E^{(i)}(x)\)은 \(\Mmod_{i-1}\), 또는 \(\Mmod^*_{i-1}\)의 가장 높은 초등인자이며, 정리 II와 제4절에 따라 \(\mideal/\gideal_{i-1}\)에 속하면서 \(x_1,\ldots,x_{i-1}\)과 무관한 모든 다항식의 다항식 의미의 최대공약원으로 정의되었다. 따라서 \(d^{(i+1)}D^{(i)}\)은 \(E^{(i)}(x)\)으로 나누어지고, \(E^{(i)}(x)\)가 \(x_i\)에 관하여 정칙이므로 \(D^{(i)}(x)\) 자체도 \(E^{(i)}(x)\)으로 나누어진다.

처음부터 변환 (12)을 (35)와 합성했다고 해도, 즉 부정원 \(u_{\mu\nu}\)를 \(w_{\mu\nu}\)로 바꾸었다고 해도 같은 논증을 수행할 수 있다. 그러면 \(H^{(i)}(z,x)\)이 대응하는 \(\bar E^{(i)}(z,x)\)으로 나누어진다. \(D^{(i)}(x),H^{(i)}(z,x)\)과 마찬가지로 \(E^{(i)}(x),\bar E^{(i)}(z,x)\), 또 \(R^{(i)}(x),\bar R^{(i)}(z,x)\)도 서로 특수화하여 얻어지므로, 일차인자 분해에서의 중복도들도 서로 일치해야 한다. 여기에 \(E^{(i)}(x)\)의 어떤 거듭제곱이 \(R^{(i)}(x)\)으로 나누어진다는 사실을 고려하면 다음을 얻는다.

\medskip
\noindent\textbf{정리 XIII.} \emph{\(R^{(i)}(x)\)을 \(P(u,x_{i+1},\ldots,x_n)\)의 한 대수적 확대체에서 일차인자 \(x_i-\bar x_i\)들로 분해하면, \(\bar x_i\)를 서로 결부된 값의 계
\[
\bar x_1,\ldots,\bar x_i,x_{i+1},\ldots,x_n
\]
로 보충하는 방법은 유일하며, 이 값의 계는 \(\mideal/\gideal_{i-1}\), 따라서 \(\mideal\)의 영점이 된다. \(R^{(i)}\)의 일차인자들에서 이렇게 생기는 \(\mideal\)의 유한개의 영점, 즉 \(x_{i+1},\ldots,x_n\)을 첨가한 뒤의 유한개의 영점은 명시적인 분해
\begin{equation}
\bar R^{(i)}(z,x)=
\prod_\nu
\left\{z_i-(\bar x_{i\nu}+v_1\bar x_{1\nu}
+\cdots+v_{i-1}\bar x_{i-1,\nu})\right\}^{\lambda_\nu}
\tag{36}
\end{equation}
에 함께 나타난다. \(\bar R^{(i)}(z,x)\)이 \(z_i\)에 관하여 정칙이므로, \(P(u)\)에 첨가한 \(x_{i+1},\ldots,x_n\)을 \(P(u)\), 또는 거기서 유도되는 대수적으로 닫힌 체의 임의의 특수한 값의 계 \(\bar x_{i+1},\ldots,\bar x_n\)으로 바꾸어도 이 분해는 유지된다.}

이로써 \(\mideal\)의 영점들이 종결식형의 개별 인자들에 따라 동시에 분리되었다. \(P(u)\)에 첨가된 부정원 \(x\)의 수는 대응하는 대수적 도형의 차원이라고도 부른다.

\(\bar R^{(i)}(z,x)\), 또는 \(R^{(i)}(x)\)의 대응 분해에서 \(P(u,x_{i+1},\ldots,x_n)\)에 관하여 켤레인 인자들을 묶자. 일반적인 \(P\)에서는 이 인자들이 전부 또는 일부 동일할 수 있다. 그러면 \(R^{(i)}(x)\)을 \(P(u,x_{i+1},\ldots,x_n)\)에 관하여 기약인 다항식들의 거듭제곱으로 분해한 식
\[
R^{(i)}(x)=S_1^{(i)}(x)^{\beta_1}\cdots
S_\nu^{(i)}(x)^{\beta_\nu}
\]
을 얻는다. 정리 X에 따라 이 분해에는 표현
\[
\gideal_i=[\mathfrak s_1,\ldots,\mathfrak s_\nu]
\]
이 대응한다. 여기서 \(S_\mu^{(i)}(x)^{\beta_\mu}\)은 \(\gideal_{i-1}\)의 아이디얼 \(\mathfrak s_\mu\)에 대한 노름, 또는 모듈 \(\Gmod_{i-1}\)의 \(\mathfrak s_\mu\)에 대응하는 모듈에 대한 노름이다. 따라서 지수 \(\beta_\mu\)의 의미도 분명해진다. 특히 \(\gamma_\mu\)가 \(S_\mu^{(i)}\)의 차수라면 중복도 \(\beta_\mu\gamma_\mu\)는 \(P(u,x_{i+1},\ldots,x_n)\)에 관하여 \(\gideal_{i-1}\)의 \(\mathfrak s_\mu\)에 대한 선형독립인 잉여류의 수와 같다.

끝으로 \(P\)의 표수가 영인 특수한 경우를 간단히 보자. \(n\)개 행렬식 \(C^{(i)}(x)\), \(i=1,2,\ldots,n\), 이 각각 \(x_i\)에 관하여 정칙으로 남도록 \(u_{\mu\nu}\)를 \(P\)의 원소 \(\bar u_{\mu\nu}\)로 특수화하면\srcfn{19)}{헨첼은 처음부터 \(u_{\mu\nu}\)에 이 특수화를 가한다. 그래서 \(H^{(i)}(z,x)\), 또는 \(\bar R^{(i)}(z,x)\)이 \(z,v\)에 관한 일차인자들로 분해됨을 보이는 데 훨씬 복잡한 고찰을 써야 한다. 이때 나타나는 지수들은 위의 지수들과 반드시 일치하지는 않는다. (E. N.)}, \(R^{(i)}\)의 정칙성도 유지된다. 이 특수화 아래 \(R^{(i)}\), 그리고 \(\bar R^{(i)}(z,x)\)은 \(\Mmod^*_{i-1}\)의 모든 \(\rho\)차 행렬식의 공통인자로 남지만 반드시 최대공약원일 필요는 없다. 그러나 이 마지막 성질도 유지되도록 언제나 특수화할 수 있다. 실제로 모듈기저의 존재가 보여 주듯 \(\bar R^{(i)}(z,x)\)은 \(\Mmod^*_{i-1}\)에서 얻는 유한개의 \(\rho\)차 행렬식의 최대공약원으로도 특징지을 수 있다. 이렇게 특수화한 \(\bar R^{(i)}(z,x)\)의 분해는 (36)에서 \(u_{\mu\nu}\)를 \(\bar u_{\mu\nu}\)로 바꾸어 단순히 얻어지므로, 소거론 전체가 유지된다.

\begin{center}
(1922년 3월 17일 접수)
\end{center}
